% !TeX root = ../../Book1.tex
% 扩充编号：A04、A05、A06。可数独立实变量、大数律与有限方差中心极限定理。
% 初稿写作与必要自检不计正式审核；文献实读范围见各节注释及来源记录。
\OptionalChapter{测度概率与极限定理}{b1:ch:30}
\BookChapterNumber{b1:ch:30}

\begin{chapterabstract}
重复观测的平均值为什么趋于稳定，平均值周围的波动又遵循什么规律？
测度积分给出随机变量与期望的语言，独立性则使不同观测之间的关系可以计算。
本章先构造可数独立随机变量列，再用概率估计、截断与极大不等式证明大数定律。
随后以特征函数描述分布，把独立和变成函数乘积，并建立中心极限定理。
两类极限定理分别解释平均值的稳定和平方根尺度上的波动。
\end{chapterabstract}

\begin{learninggoals}
\begin{itemize}
\item\label{b1:item:30-goal-model}用可测映射和推前测度表述随机变量、分布与独立性，理解可数独立变量列的构造。
\item\label{b1:item:30-goal-modes}区分几乎必然、依概率、积分及依分布收敛，并判断它们的适用对象。
\item\label{b1:item:30-goal-lln}运用 Borel--Cantelli 引理、截断和极大估计，证明可积独立同分布变量的弱、强大数律。
\item\label{b1:item:30-goal-characteristic}掌握特征函数的唯一性与连续性定理，说明紧性在极限分布存在证明中的作用。
\item\label{b1:item:30-goal-clt}证明有限方差中心极限定理，正确处理中心化、归一化、退化方差及具体分布模型。
\end{itemize}
\end{learninggoals}

本章可以在\chapref{b1:ch:12}之后阅读。
前半章需要测度的连续性、可测映射、积分收敛定理及\chapref{b1:ch:07}的有限乘积测度；
后半章使用\chapref{b1:ch:12}的弱收敛、分布函数判据和概率测度的紧性。
Gaussian 密度的归一化承接第7章的积分计算，所需特征函数公式在本章建立。

\ProseParagraphBreak

\secref{b1:sec:3001}把已有测度语言改写为随机变量的语言，
并说明比较随机变量与比较其分布的区别。
\secref{b1:sec:3002}说明独立性怎样延伸到函数及变量组，并构造给定实分布的可数乘积。
\secref{b1:sec:3003}研究事件发生无穷多次的概率。
这些工具在\secref{b1:sec:3004}至\secref{b1:sec:3005}中控制平均值：
先处理有界或有限方差变量，再借截断恢复只有一阶矩的情形。

\ProseParagraphBreak

大数律给出平均值的极限，却没有描述误差的分布。
\secref{b1:sec:3006}把分布化为特征函数，证明这种表示的唯一性及其对极限的识别能力；
\secref{b1:sec:3007}据此研究中心化独立和的平方根尺度。
阅读时要分清定理使用的是一阶矩还是二阶矩，也要记住：依分布收敛并不要求随机变量在同一个样本点上收敛。

\ProseParagraphBreak

基础概率语言及有限方差弱大数律可对照 Sheldon Axler 的
《Measure, Integration \& Real Analysis》\cite[Chapter~12]{Axler2020Measure}。
David H. Fremlin 的《Measure Theory》第二卷\cite[273, 285]{Fremlin2016Measure}
分别讨论强大数律与特征函数；本章把所需的截断和极限步骤放在使用处展开。

% !TeX root = ../../Book1.tex
\section{概率空间、随机变量与分布}
\label{b1:sec:3001}

测度论把可测集合赋予大小，概率论则把总质量归一为一，并把可测函数解释为一次试验所产生的数值。这样的解释并不改变积分和收敛定理，却提出了新的问题：同一分布能否由不同的随机变量实现，几种收敛分别记录了什么，以及平均值的稳定性需要哪些额外条件。本节先把这些问题放回已经建立的测度论框架。

\ProseParagraphBreak

基本用语可对照 Axler 的《Measure, Integration \& Real Analysis》\cite[Definition 12.1, Definition 12.13, Definition 12.18, Definition 12.27, pp.~381, 385, 388, 392]{Axler2020Measure}。本节用推前分布统一处理离散与连续情形；依分布收敛则直接承接\chapref{b1:ch:12}的测度弱收敛。

\subsection{概率空间与事件}
\label{b1:subsec:300101}

\begin{definition}\label{b1:def:30-probability-space}
设 \(\Omega\) 是集合，\(\mathcal F\) 是 \(\Omega\) 上的 \(\sigma\)-代数，\(\mathbb P\) 是 \((\Omega,\mathcal F)\) 上满足 \(\mathbb P(\Omega)=1\) 的正测度。三元组 \((\Omega,\mathcal F,\mathbb P)\) 称为概率空间，\(\Omega\) 称为\term[yangben-kongjian]{样本空间}[sample space]，\(\mathcal F\) 中的集合称为\term[shijian]{事件}[event]。对事件 \(A\)，数 \(\mathbb P(A)\) 称为它的概率。
\end{definition}

概率测度与概率空间在\secref{b1:sec:0201}已作为有限测度的例子出现。这里改用 \(\Omega,\mathcal F,\mathbb P\) 是为了适应概率论的惯用记号；可数可加性、测度连续性以及积分的含义均保持原来的定义。对可测集 \(A_n\)，例如有
\[
 \mathbb P\left(\bigcup_{n\geq1}A_n\right)
 \leq\sum_{n\geq1}\mathbb P(A_n),
 \quad
 A_n\uparrow A\ \Longrightarrow\ \mathbb P(A_n)\uparrow\mathbb P(A).
\]
若 \(A_n\downarrow A\)，上连续性也总能使用，因为 \(\mathbb P(A_1)\leq1<\infty\)。

\begin{definition}\label{b1:def:30-almost-surely}
设 \((\Omega,\mathcal F,\mathbb P)\) 是概率空间。若事件 \(A\in\mathcal F\) 满足 \(\mathbb P(A)=1\)，则称 \(A\) \term[jihu-biran]{几乎必然}[almost surely]发生，常缩写为 a.s.。一个关于 \(\omega\in\Omega\) 的性质几乎必然成立，是指它在某个可测零集之外成立。
\end{definition}

这正是前文“几乎处处”的概率论说法。可数个几乎必然成立的性质可以同时成立：若各性质的例外零集为 \(N_j\)，则 \(\mathbb P(\bigcup_jN_j)=0\)。这里的可数性不能随意删去。在 \([0,1]\) 上取 Lebesgue 概率测度，每个单点都为零测集，但全部单点的并是整个空间；因而“对每个固定参数，性质几乎必然成立”不自动给出“几乎必然对所有参数同时成立”。

\ProseParagraphBreak

有限样本空间可以取 \(\Omega=\{0,1\}\)、\(\mathcal F=\mathcal P(\Omega)\)，并指定 \(\mathbb P(\{1\})=p\)、\(\mathbb P(\{0\})=1-p\)，其中 \(0\leq p\leq1\)。连续样本空间则可取 \(([0,1],\mathcal B([0,1]),m_1)\)。后一空间中某个固定点的概率为零，点本身仍是样本空间中的可能取值；概率为零与集合为空是不同的断言。需要完备性时，可以按\subsecref{b1:subsec:020204}的办法完备化，但本章的概率空间不预先假设完备。

\subsection{随机变量与推前分布}
\label{b1:subsec:300102}

试验的结果未必本来就是一个数。把每个结果 \(\omega\) 对应到观测值 \(X(\omega)\)，才得到可用于积分和比较的函数。要求它可测，是为了使“观测值落入指定 Borel 集”确实构成事件。

\begin{definition}\label{b1:def:30-random-variable}
设 \((\Omega,\mathcal F,\mathbb P)\) 是概率空间。一个 \(\mathcal F/\mathcal B(\mathbb R)\)-可测函数 \(X:\Omega\to\mathbb R\) 称为实\term[suiji-bianliang]{随机变量}[random variable]。规定
\[
 \{X\in B\}\coloneqq\{\,\omega\in\Omega\mid X(\omega)\in B\,\},
 \quad B\in\mathcal B(\mathbb R),
\]
并在概率记号中简写为 \(\mathbb P(X\in B)\)。随机变量 \(X\) 的\term[gailv-fenbu]{概率分布}[probability distribution]是 Borel 概率测度
\[
 \mu_X\coloneqq X_\#\mathbb P,
 \quad \mu_X(B)=\mathbb P(X\in B).
\]
它的\term[fenbu-hanshu]{分布函数}[distribution function]为
\[
 F_X(t)\coloneqq\mathbb P(X\leq t)=\mu_X(({-\infty},t]),
 \quad t\in\mathbb R.
\]
\end{definition}
\symidx{\(\mu_X\)：随机变量 \(X\) 的概率分布}
\symidx{\(F_X\)：随机变量 \(X\) 的分布函数}

逆像保持不交并，故 \(\mu_X\) 的可数可加性由 \(\mathbb P\) 的可数可加性直接得到；又有 \(\mu_X(\mathbb R)=1\)。所以这里确实得到了概率测度。以后在本章中说“\(X\) 的分布”，均指这个推前测度，与\chapref{b1:ch:20}中作为测试函数连续线性泛函的分布含义不同。

\ProseParagraphBreak

测度连续性说明 \(F_X\) 单调不减、右连续，并满足
\[
 \lim_{t\to-\infty}F_X(t)=0,
 \quad \lim_{t\to+\infty}F_X(t)=1.
\]
右连续性来自 \(({-\infty},t_j]\downarrow({-\infty},t]\) 当 \(t_j\downarrow t\)。区间差给出
\[
 \mu_X((a,b])=F_X(b)-F_X(a)\quad(a<b).
\]
因此分布函数决定全部分布：半直线 \(({-\infty},t]\) 连同 \(\mathbb R\) 构成生成 \(\mathcal B(\mathbb R)\) 的 \(\pi\)-系统，有限测度唯一性适用。这也解释了分布函数为何足以描述一维概率。

\begin{definition}\label{b1:def:30-equal-laws}
设 \(X\)、\(Y\) 分别是概率空间 \((\Omega,\mathcal F,\mathbb P)\)、\((\Omega',\mathcal F',\mathbb P')\) 上的实随机变量。若它们的 Borel 分布相等，即 \(\mu_X=\mu_Y\)，则称它们\term[tong-fenbu]{同分布}[identically distributed]。对同一概率空间上的有限个实随机变量 \(X_1,\ldots,X_r\)，可测映射 \((X_1,\ldots,X_r):\Omega\to\mathbb R^r\) 的推前测度称为它们的\term[lianhe-fenbu]{联合分布}[joint distribution]。
\end{definition}

同分布不要求两个变量在同一个空间上，也不要求它们逐点相等。在 \([0,1]\) 上令 \(X(\omega)=\omega\)、\(Y(\omega)=1-\omega\)，它们的分布都是 \([0,1]\) 上的均匀概率测度，但除 \(\omega=1/2\) 外总有 \(X(\omega)\ne Y(\omega)\)。此外，单个变量的分布不决定联合分布：\((X,X)\) 总在对角线上取值，而两个独立均匀变量的联合分布将在下一节由乘积测度描述。

\ProseParagraphBreak

在前面的两点空间上，坐标变量 \(X(\omega)=\omega\) 的分布为
\[
 \mu_X=(1-p)\delta_0+p\delta_1.
\]
这个测度称为参数 \(p\) 的\term[bernoulli-fenbu]{Bernoulli 分布}[Bernoulli distribution]。任何事件 \(A\) 的示性函数 \(\mathbf1_A\) 都有这样的分布，其中 \(p=\mathbb P(A)\)。若 \(a<b\)，则密度为 \((b-a)^{-1}\cdot\mathbf1_{[a,b]}\) 的概率测度称为 \([a,b]\) 上的\term[junyun-fenbu]{均匀分布}[uniform distribution]。它们分别给出最基本的离散与连续例子。

\subsection{期望、矩与积分}
\label{b1:subsec:300103}

积分本来就用于把函数值按集合大小加权。概率测度的总质量为一，因此同一积分现在表示一次观测的平均数。

\begin{definition}\label{b1:def:30-expectation}
设 \((\Omega,\mathcal F,\mathbb P)\) 是概率空间，\(X\) 是其上的实随机变量。当 \(X\geq0\) 时，定义它的\term[qiwang]{期望}[expectation]为
\[
 \mathbb E X\coloneqq\int_\Omega X\,\mathrm d\mathbb P\in[0,\infty].
\]
同一非负积分约定也用于非负扩展实值可测函数。对一般实随机变量，若 \(\mathbb E X^+\)、\(\mathbb E X^-\) 至少有一个有限，定义 \(\mathbb E X=\mathbb E X^+-\mathbb E X^-\)；若两者都无穷，期望不定义。称 \(X\) 可积，是指 \(\mathbb E|X|<\infty\)，即 \(X\in L^1(\mathbb P)\)。对可测复值函数 \(Z\) 且 \(\mathbb E|Z|<\infty\)，规定
\[
 \mathbb E Z=\mathbb E(\operatorname{Re}Z)
 +\mathrm i\,\mathbb E(\operatorname{Im}Z).
\]
\end{definition}
\symidx{\(\mathbb E X\)：随机变量的期望}

于是 \(\mathbb E\mathbf1_A=\mathbb P(A)\)，且积分的线性性、单调收敛、Fatou 引理与控制收敛全部可以用 \(\mathbb E\) 改写。可积随机变量通常仍按几乎必然相等取等价类；谈到逐点事件时则选取可测代表。只在零测集上任意改值未必保持非完备空间上的可测性，因此不能省去代表的可测性要求。

\begin{proposition}[按分布计算期望]\label{b1:prop:30-law-integration}
设 \(X\) 是概率空间 \((\Omega,\mathcal F,\mathbb P)\) 上的实随机变量，分布为 \(\mu_X\)。若 \(h:\mathbb R\to[0,\infty]\) 是 Borel 函数，则
\[
 \mathbb E\bigl(h(X)\bigr)=\int_{\mathbb R}h(x)\,\mathrm d\mu_X(x).
\]
若 \(h:\mathbb R\to\mathbb C\) 为 Borel 函数，则 \(h(X)\in L^1(\mathbb P)\) 当且仅当 \(h\in L^1(\mu_X)\)，并且可积时同一等式成立。
\end{proposition}
\begin{proof}
当 \(h=\mathbf1_B\)、\(B\in\mathcal B(\mathbb R)\) 时，等式就是 \(\mu_X(B)=\mathbb P(X\in B)\)。有限非负线性组合给出非负简单函数情形，再以非负简单函数递增逼近 \(h\)，使用\cref{b1:thm:monotone-convergence}。对复值函数先把结论用于 \(|h|\)，得到可积性的等价；再拆分实部、虚部的正负部分，即得所需等式。
\end{proof}

该公式说明，期望和绝对矩只依赖分布。同一证明也适用于 \(\mathbb R^r\) 值可测映射与 Borel 函数 \(h:\mathbb R^r\to\mathbb C\)。若 \(\mu_X\ll m_1\)，则其 Radon--Nikodym 密度 \(f\geq0\) 称为 \(X\) 的\term[gailv-midu]{概率密度}[probability density]，满足 \(\int_{\mathbb R}f\,\mathrm dx=1\)，并且
\[
 \mathbb E\bigl(h(X)\bigr)=\int_{\mathbb R}h(x)\cdot f(x)\,\mathrm dx
\]
对非负或可积的 \(h\) 成立。分布不一定有密度，例如 Bernoulli 分布集中在两个点上。

\begin{definition}\label{b1:def:30-moments-variance}
设 \(X\) 是概率空间 \((\Omega,\mathcal F,\mathbb P)\) 上的实随机变量。对 \(p>0\)，\(\mathbb E|X|^p\) 称为 \(p\) 阶\term[jue-dui-ju]{绝对矩}[absolute moment]；当整数 \(k\geq1\) 且 \(\mathbb E|X|^k<\infty\) 时，\(\mathbb E(X^k)\) 称为 \(k\) 阶\term[ju]{矩}[moment]。若 \(X\in L^2(\mathbb P)\)，则定义它的\term[fangcha]{方差}[variance]与\term[biaozhun-cha]{标准差}[standard deviation]为
\[
 \operatorname{Var}(X)\coloneqq\mathbb E\left((X-\mathbb E X)^2\right),
 \quad \sigma_X\coloneqq\sqrt{\operatorname{Var}(X)}.
\]
\end{definition}
\symidx{\(\operatorname{Var}(X)\)：实随机变量的方差}

\begin{proposition}\label{b1:prop:30-variance-formula}
设 \(X\in L^2(\mathbb P)\) 是概率空间 \((\Omega,\mathcal F,\mathbb P)\) 上的实随机变量。则 \(X\in L^1(\mathbb P)\)，并且
\[
 \operatorname{Var}(X)=\mathbb E(X^2)-(\mathbb E X)^2.
\]
对每个 \(a\in\mathbb R\)，还有
\[
 \mathbb E\left((X-a)^2\right)
 =\operatorname{Var}(X)+(\mathbb E X-a)^2.
\]
特别地，\(\operatorname{Var}(X)=0\) 当且仅当 \(X=\mathbb E X\) 几乎必然。
\end{proposition}
\begin{proof}
由 Cauchy--Schwarz 不等式及 \(\mathbb P(\Omega)=1\)，有 \(\mathbb E|X|\leq\|X\|_2\)。将 \((X-\mathbb E X)^2\) 展开并利用期望的线性性，得到第一式。再写
\[
 X-a=(X-\mathbb E X)+(\mathbb E X-a).
\]
平方后交叉项的期望为零，因为 \(\mathbb E(X-\mathbb E X)=0\)，故得到第二式。最后由\cref{b1:prop:integral-markov}的非负零积分判别，\(\mathbb E((X-\mathbb E X)^2)=0\) 等价于 \(X=\mathbb E X\) 几乎必然。
\end{proof}

这表明期望是使均方误差最小的常数，方差则是这个最小误差。对参数 \(p\) 的 Bernoulli 变量，\(X^2=X\)，所以
\[
 \mathbb E X=p,
 \quad \operatorname{Var}(X)=p(1-p).
\]
对 \([a,b]\) 上的均匀变量，直接积分得到
\[
 \mathbb E X=\frac{a+b}{2},
 \quad \operatorname{Var}(X)=\frac{(b-a)^2}{12}.
\]
这些计算只用到了分布，还没有涉及不同变量之间的依赖关系。

\subsection{随机变量的收敛方式}
\label{b1:subsec:300104}

随机变量列既可以逐点比较，也可以只比较误差事件的概率，或者只比较各自的分布。几乎必然收敛记录了单个样本的长期行为；依概率收敛只要求每个固定精度下的大误差越来越少；依分布收敛则允许完全不比较同一试验中的取值。

\begin{definition}\label{b1:def:30-convergence-modes}
设 \(X_n,X\) 是同一概率空间 \((\Omega,\mathcal F,\mathbb P)\) 上的实随机变量。
\begin{enumerate}
 \item 若 \(\mathbb P(\{\omega:X_n(\omega)\to X(\omega)\})=1\)，称 \(X_n\) 几乎必然收敛于 \(X\)，记为 \(X_n\xrightarrow{\mathrm{a.s.}}X\)。
 \item 若对每个 \(\varepsilon>0\)，有 \(\mathbb P(|X_n-X|>\varepsilon)\to0\)，称 \(X_n\) \term[yi-gailv-shoulian]{依概率收敛}[convergence in probability]于 \(X\)，记为 \(X_n\xrightarrow{\mathbb P}X\)。
 \item 对 \(1\leq p\leq\infty\)，若 \(X_n,X\in L^p(\mathbb P)\)，且 \(\|X_n-X\|_p\to0\)，称 \(X_n\) 在 \(L^p\) 中收敛于 \(X\)。
 \item 若 Borel 分布满足 \(\mu_{X_n}\Rightarrow\mu_X\)，称 \(X_n\) 依分布收敛于 \(X\)，记为 \(X_n\xrightarrow{\mathrm d}X\)。等价地，对每个 \(h\in C_b(\mathbb R)\)，有 \(\mathbb E(h(X_n))\to\mathbb E(h(X))\)。这一项也适用于分处不同概率空间的随机变量。
\end{enumerate}
\end{definition}

依概率收敛就是\cref{b1:def:convergence-in-measure}在总质量为一的情形；依分布收敛已经在\subsecref{b1:subsec:120104}引入。这里明确列出它们，是为了避免把“误差变小”与“分布接近”当作同一件事。

\begin{proposition}\label{b1:prop:30-convergence-implications}
设 \(X_n,X\) 是概率空间 \((\Omega,\mathcal F,\mathbb P)\) 上的实随机变量。则：
\begin{enumerate}
 \item 几乎必然收敛推出依概率收敛；对 \(1\leq p\leq\infty\)，\(L^p\) 收敛也推出依概率收敛。
 \item \(X_n\xrightarrow{\mathbb P}X\) 当且仅当任意子列都有一个进一步的子列几乎必然收敛于 \(X\)。
 \item 依概率收敛推出依分布收敛。
 \item 若 \(c\in\mathbb R\) 为常数，则 \(X_n\xrightarrow{\mathrm d}c\) 当且仅当 \(X_n\xrightarrow{\mathbb P}c\)。
\end{enumerate}
\end{proposition}
\begin{proof}
第一项的几乎必然部分和第二项的正向结论，均由\cref{b1:lem:measure-convergence-subsequence}得到：它的从几乎处处到依测度的方向要求有限总质量，这里恰有 \(\mathbb P(\Omega)=1\)。对 \(p<\infty\)，\cref{b1:prop:integral-markov}给出
\[
 \mathbb P(|X_n-X|>\varepsilon)
 \leq\varepsilon^{-p}\cdot\mathbb E|X_n-X|^p
 =\varepsilon^{-p}\cdot\|X_n-X\|_p^p.
\]
对 \(p=\infty\)，充分大的 \(n\) 满足 \(\|X_n-X\|_\infty<\varepsilon\)，所需事件为零测集。

若第二项的逆向不成立，则某个 \(\varepsilon>0\)、\(\delta>0\) 与子列 \((X_{n_k})\) 满足 \(\mathbb P(|X_{n_k}-X|>\varepsilon)\geq\delta\) 对每个 \(k\) 成立。该子列不可能有几乎必然收敛于 \(X\) 的进一步子列，否则第一项使这些概率趋零，矛盾。

对第三项，固定 \(h\in C_b(\mathbb R)\)。依概率收敛的任意子列都可再抽出几乎必然收敛于 \(X\) 的子列。在后者上，连续性和控制收敛给出 \(\mathbb E(h(X_{n_{k_j}}))\to\mathbb E(h(X))\)，控制函数为常数 \(\|h\|_\infty\)。于是原数列的测试期望也必收敛到此值；否则可抽出与此值保持正距离的子列。由定义得到依分布收敛。

最后，若 \(\mu_{X_n}\Rightarrow\delta_c\)，对每个 \(\varepsilon>0\)，闭集 \(F_\varepsilon=\{x:|x-c|\geq\varepsilon\}\) 满足 \(\delta_c(F_\varepsilon)=0\)。\cref{b1:thm:portmanteau}的闭集判据给出
\[
 \limsup_n\mathbb P(|X_n-c|\geq\varepsilon)\leq0.
\]
因此依概率收敛成立，反向则是第三项。
\end{proof}

这些蕴含一般不能倒转。在 \(([0,1),\mathcal B([0,1)),m_1)\) 上，按每一层二进区间依次定义
\[
 X_{2^m+k}=\mathbf1_{[k2^{-m},(k+1)2^{-m})},
 \quad m\geq0,\quad 0\leq k<2^m.
\]
当指标处在第 \(m\) 层时，\(\mathbb P(X_n=1)=2^{-m}\)，所以 \(X_n\) 依概率且在每个有限 \(L^p\)、\(p\geq1\) 中收敛于零。但每个 \(\omega\in[0,1)\) 在每层恰好落入一个区间，故 \(X_n(\omega)=1\) 出现无穷多次；每层的其余区间又使零值出现无穷多次。这个序列在任何样本点都不收敛。

\ProseParagraphBreak

反过来，在 \((0,1)\) 的 Lebesgue 概率空间上取 \(Y_n=n\cdot\mathbf1_{(0,1/n)}\)。每个固定 \(\omega>0\) 最终离开 \((0,1/n)\)，所以 \(Y_n\to0\) 处处成立，但 \(\mathbb E|Y_n|=1\)，没有 \(L^1\) 收敛。小概率事件上的大取值保留了全部积分，这正是\secref{b1:sec:0505}中一致可积条件所控制的现象。

\ProseParagraphBreak

依分布收敛更弱。在等概率两点空间上令 \(X(\omega)=\omega\)，并令 \(Z_{2n}=X\)、\(Z_{2n-1}=1-X\)。全部 \(Z_n\) 都有参数 \(1/2\) 的 Bernoulli 分布，所以它们依分布收敛于 \(X\)；但 \(\mathbb P(|Z_{2n-1}-X|>1/2)=1\)，因而不依概率收敛于 \(X\)。实际上相邻两项处处相差一，所以这个序列也不可能依概率收敛于其他变量。后面的大数定律与中心极限定理将分别使用这些不同的收敛目标。

% !TeX root = ../../Book1.tex
\section{独立性与可数乘积概率空间}
\label{b1:sec:3002}

只知道每个变量各自的分布，还不能确定它们的和或乘积如何变化。独立性给出一种精确的联合结构：对不同变量分别提出的条件，其同时成立的概率等于各自概率的乘积。先在事件上写出的乘法关系，必须能够推广到生成的 \(\sigma\)-代数，才能用于变量的函数、分块求和以及极限。

\ProseParagraphBreak

独立性的三种表达可对照 Fremlin 的《Measure Theory》第二卷\cite[Definition 272A, Proposition 272D, Theorem 272G, Proposition 272K]{Fremlin2016Measure}。本节先用\cref{b1:thm:pi-lambda}完成这些推广，再为给定的一列实 Borel 分布构造乘积概率空间；后一步将实际提供后文所需的无限独立序列。

\subsection{事件与子代数的独立性}
\label{b1:subsec:300201}

\begin{definition}\label{b1:def:30-independent-events}
设 \((\Omega,\mathcal F,\mathbb P)\) 是概率空间，\((A_i)_{i\in I}\) 是事件族。若对任意互异指标 \(i_1,\ldots,i_r\in I\)、\(r\geq1\)，都有
\[
 \mathbb P\left(\bigcap_{j=1}^rA_{i_j}\right)
 =\prod_{j=1}^r\mathbb P(A_{i_j}),
\]
则称这族事件\term[duli-shijian]{相互独立}[mutually independent events]，也简称独立。若只要求这个等式在 \(r=2\) 时成立，则称它们\term[liangliang-duli]{两两独立}[pairwise independent]。
\end{definition}

定义始终检查有限子族，即使指标集 \(I\) 是无限的也如此。只检查成对关系一般不够。在等概率空间 \(\Omega=\{0,1\}^2\) 上令
\[
 A=\{(u,v):u=0\},
 \quad B=\{(u,v):v=0\},
 \quad C=\{(u,v):u=v\}.
\]
三个事件各有概率 \(1/2\)，任意两个的交有概率 \(1/4\)，但
\[
 \mathbb P(A\cap B\cap C)=\frac14\ne\frac18
 =\mathbb P(A)\cdot\mathbb P(B)\cdot\mathbb P(C).
\]
因此两两独立并不等于相互独立。

\begin{definition}\label{b1:def:30-independent-sigma}
设 \((\Omega,\mathcal F,\mathbb P)\) 是概率空间，\((\mathcal F_i)_{i\in I}\) 是 \(\mathcal F\) 的子 \(\sigma\)-代数族。若对任意互异指标 \(i_1,\ldots,i_r\) 及任意 \(A_j\in\mathcal F_{i_j}\)，均有
\[
 \mathbb P\left(\bigcap_{j=1}^rA_j\right)
 =\prod_{j=1}^r\mathbb P(A_j),
\]
则称 \((\mathcal F_i)_{i\in I}\) 是独立的 \(\sigma\)-代数族。
\end{definition}

每个 \(\mathcal F_i\) 可以看作关于第 \(i\) 个观测所能提出的全部可测条件。实际核验时往往只知道一个生成族上的乘法关系。下面的结论说明，只要生成族对有限交封闭，这些信息已经足够。

\begin{proposition}[独立性的生成族判据]\label{b1:prop:30-independent-pi-lambda}
设 \((\Omega,\mathcal F,\mathbb P)\) 是概率空间。对每个 \(i\in I\)，设 \(\mathcal P_i\subseteq\mathcal F\) 是包含 \(\Omega\) 的 \(\pi\)-系统，令 \(\mathcal F_i=\sigma(\mathcal P_i)\)。若对任意互异的 \(i_1,\ldots,i_r\) 及 \(A_j\in\mathcal P_{i_j}\)，都有
\[
 \mathbb P\left(\bigcap_{j=1}^rA_j\right)
 =\prod_{j=1}^r\mathbb P(A_j),
\]
则 \((\mathcal F_i)_{i\in I}\) 是独立的。
\end{proposition}
\begin{proof}
固定互异指标 \(i_1,\ldots,i_r\)。先固定 \(A_j\in\mathcal P_{i_j}\)、\(2\leq j\leq r\)，并记
\[
 H=\bigcap_{j=2}^rA_j,
 \quad c=\prod_{j=2}^r\mathbb P(A_j),
 \quad \mathcal D=\{A\in\mathcal F_{i_1}:\mathbb P(A\cap H)=c\cdot\mathbb P(A)\}.
\]
空交与空积分别取 \(\Omega\) 与一。原假设给出 \(\mathbb P(H)=c\)，所以 \(\Omega\in\mathcal D\)。若 \(A\subseteq B\) 且 \(A,B\in\mathcal D\)，有限测度的差公式给出
\[
 \mathbb P((B\setminus A)\cap H)
 =c\cdot\bigl(\mathbb P(B)-\mathbb P(A)\bigr)
 =c\cdot\mathbb P(B\setminus A).
\]
对两两不交的 \(A_n\in\mathcal D\)，可数可加性又给出它们的并属于 \(\mathcal D\)。因此 \(\mathcal D\) 是 \(\lambda\)-系统。它包含 \(\mathcal P_{i_1}\)，故由\cref{b1:thm:pi-lambda}包含 \(\mathcal F_{i_1}\)。这就把第一个坐标的集合由生成族推广到了整个 \(\sigma\)-代数。

现在固定任意 \(A_1\in\mathcal F_{i_1}\) 及其余尚未推广的生成族集合，对第二个坐标重复同一论证。检验 \(\Omega\) 属于相应 \(\lambda\)-系统时，只需把第二个集合取为 \(\Omega\)，这是上一步已经成立的乘法等式。依次推广全部 \(r\) 个坐标，便得到任意 \(A_j\in\mathcal F_{i_j}\) 时的等式。有限指标组任意，结论成立。
\end{proof}

\begin{corollary}\label{b1:cor:30-events-generated-sigma}
设 \((A_i)_{i\in I}\) 是概率空间 \((\Omega,\mathcal F,\mathbb P)\) 上的事件族。则 \((A_i)_{i\in I}\) 独立，当且仅当 \((\sigma(A_i))_{i\in I}\) 独立。特别地，若 \((A_i)\) 独立，则任意选择 \(B_i=A_i\) 或 \(B_i=\Omega\setminus A_i\) 所得到的事件族也独立。
\end{corollary}
\begin{proof}
取 \(\mathcal P_i=\{\Omega,A_i\}\)。它是包含全集的 \(\pi\)-系统，而对其中集合的乘法等式就是原事件族对相应子族的独立性。因此前一命题适用。逆向只需在独立 \(\sigma\)-代数中选取 \(A_i\)。最后 \(A_i\) 与其补集均属于 \(\sigma(A_i)\)，所以结论也覆盖任意补集选择。
\end{proof}

\subsection{随机变量及函数族的独立性}
\label{b1:subsec:300202}

\begin{definition}\label{b1:def:30-independent-variables}
设 \((X_i)_{i\in I}\) 是概率空间 \((\Omega,\mathcal F,\mathbb P)\) 上的实随机变量族。定义
\[
 \sigma(X_i)\coloneqq\{X_i^{-1}(B):B\in\mathcal B(\mathbb R)\}.
\]
若 \((\sigma(X_i))_{i\in I}\) 独立，则称 \((X_i)_{i\in I}\) 为\term[duli-suiji-bianliang]{独立随机变量}[independent random variables]。等价地，对任意互异指标 \(i_1,\ldots,i_r\) 及 Borel 集 \(B_1,\ldots,B_r\subseteq\mathbb R\)，有
\[
 \mathbb P(X_{i_1}\in B_1,\ldots,X_{i_r}\in B_r)
 =\prod_{j=1}^r\mathbb P(X_{i_j}\in B_j).
\]
对实值或复值可测函数的独立性，也按它们所生成的 \(\sigma\)-代数定义，其中复值域采用 \(\mathcal B(\mathbb C)\)。
\end{definition}

逆像保持补集与可数并，所以 \(\sigma(X_i)\) 确实是 \(\sigma\)-代数；它也是使 \(X_i\) 可测的最小 \(\sigma\)-代数。定义中的联合概率是相应事件交集的简写。由生成族判据，只需检验每个 \(B_j\) 为半直线 \(({-\infty},t_j]\) 或 \(\mathbb R\) 的情形，就能得到全部 Borel 集的等式。

\begin{proposition}\label{b1:prop:30-independent-functions}
设 \((X_i)_{i\in I}\) 是概率空间 \((\Omega,\mathcal F,\mathbb P)\) 上的独立实随机变量族。对每个 \(i\)，设 \(h_i:\mathbb R\to\mathbb R\) 或 \(h_i:\mathbb R\to\mathbb C\) 是 Borel 函数。则可测函数族 \((h_i(X_i))_{i\in I}\) 仍独立。
\end{proposition}
\begin{proof}
对相应值域中的 Borel 集 \(B\)，有
\[
 \{h_i(X_i)\in B\}=\{X_i\in h_i^{-1}(B)\}\in\sigma(X_i).
\]
所以 \(\sigma(h_i(X_i))\subseteq\sigma(X_i)\)。独立 \(\sigma\)-代数族的任意子 \(\sigma\)-代数族仍满足定义中的乘法等式，故结论成立。
\end{proof}

\begin{proposition}[不交分块保持独立性]\label{b1:prop:30-independent-blocks}
设 \((\mathcal F_i)_{i\in I}\) 是概率空间 \((\Omega,\mathcal F,\mathbb P)\) 上的独立子 \(\sigma\)-代数族，\((J_s)_{s\in S}\) 是 \(I\) 中两两不交的子集族。令
\[
 \mathcal H_s=\sigma\left(\bigcup_{i\in J_s}\mathcal F_i\right),
 \quad s\in S.
\]
则 \((\mathcal H_s)_{s\in S}\) 独立。若可测函数 \(Y_s\) 对 \(\mathcal H_s\) 可测，则 \((Y_s)_{s\in S}\) 也独立。
\end{proposition}
\begin{proof}
对每个 \(s\)，取生成 \(\mathcal H_s\) 的 \(\pi\)-系统
\[
 \mathcal P_s=\left\{\bigcap_{i\in F}A_i:
 F\subseteq J_s\text{ 有限},\ A_i\in\mathcal F_i\right\}.
\]
其中 \(F=\varnothing\) 时取全集。它对有限交封闭，因为相同指标的集合可以先在 \(\mathcal F_i\) 内求交。若 \(s_1,\ldots,s_r\) 互异，\(C_j\in\mathcal P_{s_j}\)，各 \(C_j\) 所使用的指标组互不相交。原 \(\sigma\)-代数族的独立性因此同时给出
\[
 \mathbb P\left(\bigcap_{j=1}^rC_j\right)
 =\prod_{j=1}^r\mathbb P(C_j).
\]
这里先把左端展开为全部所用指标的有限交，再把相应概率因子按块重新组合。\cref{b1:prop:30-independent-pi-lambda}便给出 \((\mathcal H_s)\) 的独立性。又有 \(\sigma(Y_s)\subseteq\mathcal H_s\)，故最后一个断言成立。
\end{proof}

例如，若 \(X_1,X_2,\ldots\) 独立，则 \(X_1+\cdots+X_m\) 与 \(X_{m+1}+\cdots+X_n\) 在 \(m<n\) 时独立。更一般地，对不交的有限指标块分别施加任意 Borel 函数，所得到的变量仍独立。这个结论不能只凭“不同项互不影响”的直观代替，因为后文的最大值事件会同时涉及一整块变量。

\begin{proposition}[联合分布与乘积]\label{b1:prop:30-independent-joint-law}
设 \(X_1,\ldots,X_r\) 是概率空间 \((\Omega,\mathcal F,\mathbb P)\) 上的实随机变量，各自的 Borel 分布为 \(\mu_1,\ldots,\mu_r\)。则 \(X_1,\ldots,X_r\) 独立，当且仅当其联合分布等于有限乘积测度 \(\mu_1\otimes\cdots\otimes\mu_r\)。
\end{proposition}
\begin{proof}
独立性恰好说明联合分布与上述乘积测度在每个 Borel 矩形 \(B_1\times\cdots\times B_r\) 上相等。两者都是 \(\mathbb R^r\) 上的概率测度，Borel 矩形组成包含全集的生成 \(\pi\)-系统，故\cref{b1:thm:pi-lambda}给出两测度相等。反过来，若联合分布为乘积，矩形上的等式给出独立性；将部分 \(B_j\) 取为 \(\mathbb R\)，同时得到所有子族的等式。
\end{proof}

独立性还把乘积的积分拆成单变量积分。下面给出可积实、复函数都能使用的版本；其中复值情形将在特征函数的计算中出现。Axler 的《Measure, Integration \& Real Analysis》\cite[12.16--12.17, pp.~387--388]{Axler2020Measure}讨论了函数保持独立与平方可积变量的乘积期望；这里利用独立性把可积条件放宽到各个因子分别属于 \(L^1\)。

\begin{proposition}\label{b1:prop:30-independent-expectations}
设 \(Z_1,\ldots,Z_r\) 是概率空间 \((\Omega,\mathcal F,\mathbb P)\) 上的独立可测函数。
\begin{enumerate}
 \item 若各 \(Z_j\) 为非负实值函数，则
 \[
  \mathbb E\left(\prod_{j=1}^rZ_j\right)=\prod_{j=1}^r\mathbb E Z_j,
 \]
 其中允许无穷值，右端按 \(0\cdot\infty=0\) 解释。
 \item 若各 \(Z_j\) 为可积的实值或复值函数，则 \(\prod_{j=1}^rZ_j\) 可积，并且
 \[
  \mathbb E\left|\prod_{j=1}^rZ_j\right|=\prod_{j=1}^r\mathbb E|Z_j|,
  \quad
  \mathbb E\left(\prod_{j=1}^rZ_j\right)=\prod_{j=1}^r\mathbb E Z_j.
 \]
\end{enumerate}
\end{proposition}
\begin{proof}
先设各 \(Z_j\) 是简单函数，写成 \(Z_j=\sum_{k=1}^{N_j}a_{jk}\cdot\mathbf1_{A_{jk}}\)，其中 \(A_{jk}\in\sigma(Z_j)\)。展开有限乘积，并用独立性计算每个交集的概率，得到
\[
 \mathbb E\left(\prod_{j=1}^rZ_j\right)
 =\sum_{k_1,\ldots,k_r}\prod_{j=1}^r
 \bigl(a_{j k_j}\cdot\mathbb P(A_{j k_j})\bigr)
 =\prod_{j=1}^r\mathbb E Z_j.
\]
这也适用于复简单函数。

对非负 \(Z_j\)，选取对 \(\sigma(Z_j)\) 可测的非负简单函数 \(s_{jn}\uparrow Z_j\)。简单函数情形与单调收敛给出第一项。若某个 \(\mathbb E Z_j=0\)，则该 \(Z_j=0\) 几乎必然，左端也为零；否则各因子期望都严格为正，有限个递增非负数列的乘积收敛到其扩展实数乘积。这同时验证了无穷值的约定。

对可积的实、复函数，\(|Z_j|\) 仍独立，第一项给出
\[
 \mathbb E\left(\prod_{j=1}^r|Z_j|\right)
 =\prod_{j=1}^r\mathbb E|Z_j|<\infty.
\]
再选对 \(\sigma(Z_j)\) 可测的有限值简单函数 \(s_{jn}\to Z_j\)，使 \(|s_{jn}|\leq|Z_j|\)。例如实函数分别截断、二进逼近正负部分；复函数分别逼近实虚部并使各分量的绝对值不增。乘积由可积函数 \(\prod_j|Z_j|\) 控制，各因子由 \(|Z_j|\) 控制。简单函数等式两端使用控制收敛，便得到第二项。
\end{proof}

\subsection{可数乘积概率测度}
\label{b1:subsec:300203}

有限乘积测度已经由\cref{b1:thm:product-measure-construction}与\cref{b1:prop:finite-product-associativity}建立。但要讨论无限次独立试验，仍须在同一个空间上容纳全部坐标。只给出每个有限维矩形的概率还不够：必须证明这些规定具有可数可加性。

\ProseParagraphBreak

下面只构造后文实际需要的情形：各坐标取值于 \(\mathbb R\)，指定的是 Borel 概率分布。实空间有限维概率测度的紧内正则性，使我们能够用一个对角选取补上可数可加性的关键一步。更一般的概率空间乘积可参看 Fremlin 的《Measure Theory》第二卷\cite[Definition 254C, Theorem 254F]{Fremlin2016Measure}；该书把乘积取为完备测度，本节先在坐标生成的 \(\sigma\)-代数上构造唯一测度。

\begin{theorem}[实 Borel 概率的可数乘积]\label{b1:thm:30-countable-product}
设 \((\mu_j)_{j\geq1}\) 是 \((\mathbb R,\mathcal B(\mathbb R))\) 上的一列概率测度，令 \(\Omega=\mathbb R^{\mathbb N}\)，\(\xi_j(\omega)=\omega_j\)，并令
\[
 \mathcal F=\sigma\{\xi_j^{-1}(B):j\geq1,\ B\in\mathcal B(\mathbb R)\}.
\]
则 \(\mathcal F\) 上存在唯一概率测度 \(\mathbb P\)，使对每个 \(n\geq1\) 及 \(B\in\mathcal B(\mathbb R^n)\)，有
\[
 \mathbb P\bigl((\xi_1,\ldots,\xi_n)\in B\bigr)
 =(\mu_1\otimes\cdots\otimes\mu_n)(B).
\]
特别地，对互异的有限指标 \(j_1,\ldots,j_r\) 与 Borel 集 \(B_1,\ldots,B_r\)，有
\[
 \mathbb P(\xi_{j_1}\in B_1,\ldots,\xi_{j_r}\in B_r)
 =\prod_{k=1}^r\mu_{j_k}(B_k).
\]
此外，\(\mathcal F\) 正是可数乘积拓扑下 \(\mathbb R^{\mathbb N}\) 的 Borel \(\sigma\)-代数。
\end{theorem}
\begin{proof}
记 \(\pi_n=(\xi_1,\ldots,\xi_n)\)，\(q_n=\mu_1\otimes\cdots\otimes\mu_n\)。有限乘积的结合律及每个因子的总质量为一，给出一致性关系
\[
 q_m(B\times\mathbb R^{m-n})=q_n(B)
 \quad(m\geq n,\ B\in\mathcal B(\mathbb R^n)).
\]
取柱集族
\[
 \mathcal A=\{\pi_n^{-1}(B):n\geq1,\ B\in\mathcal B(\mathbb R^n)\}.
\]
把有限多个柱集都提升到共同的前 \(m\) 个坐标后，补集、交集与并集仍为同类柱集，因此 \(\mathcal A\) 是集合代数。它生成 \(\mathcal F\)：一方面坐标条件都属于 \(\mathcal A\)，另一方面 \(\pi_n\) 对 \(\mathcal F/\mathcal B(\mathbb R^n)\) 可测，所以每个柱集都属于 \(\mathcal F\)。

在 \(\mathcal A\) 上规定
\[
 p(\pi_n^{-1}(B))=q_n(B).
\]
这一定义与柱集的表示无关。事实上，若 \(m\geq n\) 且 \(\pi_n^{-1}(B)=\pi_m^{-1}(D)\)，由于任意 \(\mathbb R^m\) 中的点都能补上零坐标而延长为 \(\Omega\) 中的点，必有 \(D=B\times\mathbb R^{m-n}\)；再用一致性即可。不同表示均可先提升到共同维数，故一般情形也成立。相同办法把有限个不交柱集提升到共同维数，\(q_m\) 的有限可加性便给出 \(p\) 的有限可加性，且 \(p(\Omega)=1\)。

关键是证明：若 \(C_m\in\mathcal A\)、\(C_m\downarrow\varnothing\)，则 \(p(C_m)\to0\)。反设存在 \(\varepsilon>0\)，使全部 \(p(C_m)\geq\varepsilon\)。递归增大表示维数，写成
\[
 C_m=\pi_{k_m}^{-1}(B_m),
 \quad k_1<k_2<\cdots,\quad k_m\geq m,
 \quad B_m\in\mathcal B(\mathbb R^{k_m}).
\]
由\cref{b1:lem:polish-probability-radon}，有限维概率测度 \(q_{k_m}\) 对每个 Borel 集都紧内正则。因此可选紧集 \(K_m\subseteq B_m\)，使
\[
 q_{k_m}(B_m\setminus K_m)<\varepsilon2^{-m-2}.
\]
令 \(D_m=\pi_{k_m}^{-1}(K_m)\subseteq C_m\)。对每个 \(M\)，有限可加性与次可加估计给出
\[
 \begin{aligned}
 p\left(C_M\setminus\bigcap_{m=1}^MD_m\right)
 &\leq\sum_{m=1}^Mp(C_m\setminus D_m)\\
 &<\varepsilon\sum_{m=1}^M2^{-m-2}<\frac\varepsilon2.
 \end{aligned}
\]
因此 \(p(\bigcap_{m=1}^MD_m)>\varepsilon/2\)，特别地这个有限交非空。选取 \(x^{(M)}\in\bigcap_{m=1}^MD_m\)。

固定 \(m\) 时，所有 \(M\geq m\) 的前 \(k_m\) 个坐标都位于紧集 \(K_m\)。先从 \((x^{(M)})\) 中抽取前 \(k_1\) 个坐标收敛的子列，再从中抽取前 \(k_2\) 个坐标收敛的子列，依次继续。取这些嵌套子列的对角子列，可使每个固定前缀都收敛。各前缀极限在重叠坐标上一致；又因 \(k_m\geq m\)，它们共同确定一个 \(x\in\mathbb R^{\mathbb N}\)。每个 \(K_m\) 闭，所以 \(\pi_{k_m}(x)\in K_m\)。于是
\[
 x\in\bigcap_{m\geq1}D_m\subseteq\bigcap_{m\geq1}C_m,
\]
与交集为空矛盾。所需的空集处上连续性得证。

现设 \(E_j\in\mathcal A\) 两两不交，且 \(E=\bigcup_{j\geq1}E_j\in\mathcal A\)。余集
\[
 R_N=E\setminus\bigcup_{j=1}^NE_j
\]
属于 \(\mathcal A\)，并递减到空集。由有限可加性和刚证的连续性，
\[
 p(E)=\sum_{j=1}^Np(E_j)+p(R_N)
 \longrightarrow\sum_{j=1}^{\infty}p(E_j).
\]
故 \(p\) 是前测度。\cref{b1:thm:caratheodory-extension}将它延拓到 \(\sigma(\mathcal A)=\mathcal F\) 上；总质量仍为一，而\cref{b1:thm:sigma-finite-uniqueness}给出唯一性。有限坐标的矩形值已经包含在构造中；将未出现的坐标条件取为 \(\mathbb R\)，便得到任意有限指标组的乘法公式。

最后，每个坐标映射在乘积拓扑下连续，所以 \(\mathcal F\) 包含于相应 Borel \(\sigma\)-代数。反过来，有限多个坐标落在有理端点开区间内所给出的柱集构成可数拓扑基，而且都属于 \(\mathcal F\)。每个开集是其中某些基元的可数并，故也属于 \(\mathcal F\)。这证明二者相等。
\end{proof}

证明中有限可加性只用了有限维一致性，可数可加性则另外使用了紧内逼近和对角选取。把柱集上的乘法规定直接称为概率测度，会跳过这一步。若需要完备概率空间，可在构造完成后再对 \(\mathbb P\) 完备化；坐标变量的分布与独立性均不因此改变。

\subsection{独立随机变量列的构造}
\label{b1:subsec:300204}

\begin{corollary}\label{b1:cor:30-iid-construction}
设 \((\mu_j)_{j\geq1}\) 是实轴上的一列 Borel 概率测度。则存在概率空间 \((\Omega,\mathcal F,\mathbb P)\) 及独立实随机变量列 \((X_j)_{j\geq1}\)，使每个 \(X_j\) 的分布为 \(\mu_j\)。若全部 \(\mu_j\) 等于同一分布 \(\mu\)，则这些变量独立且同分布，简称\term[duli-tongfenbu]{独立同分布}[independent and identically distributed]，缩写为 i.i.d.。
\end{corollary}
\begin{proof}
取前一定理构造的 \(\Omega=\mathbb R^{\mathbb N}\)、\(\mathcal F\)、\(\mathbb P\)，并令 \(X_j=\xi_j\)。坐标可测性由 \(\mathcal F\) 的定义保证；有限坐标乘法公式同时给出每个边缘分布和全部有限子族的独立性。若各边缘相同，定义即给出独立同分布。
\end{proof}

例如，取每个 \(\mu_j=(1-p)\delta_0+p\delta_1\)，就得到无限 Bernoulli 试验。各坐标属于 \(\{0,1\}\) 的事件均有概率一，可数交也有概率一；所以可把空间限制在 \(\{0,1\}^{\mathbb N}\) 上。此时对 \(a_j\in\{0,1\}\)，
\[
 \mathbb P(X_1=a_1,\ldots,X_n=a_n)
 =\prod_{j=1}^n\mu_j(\{a_j\}).
\]
如果 \(0<p<1\)，右端也可写成 \(p^{\sum_ja_j}(1-p)^{n-\sum_ja_j}\)。这一构造允许讨论无限次抛掷的事件，例如某种结果是否出现无穷多次；它们已经超出单个有限维试验的定义域。

\ProseParagraphBreak

取各 \(\mu_j\) 为 \([0,1]\) 上的均匀分布，则得到独立均匀变量列。对 \(0\leq t_j\leq1\)，有
\[
 \mathbb P(X_1\leq t_1,\ldots,X_n\leq t_n)=\prod_{j=1}^nt_j.
\]
这里先有完整的概率空间，再由坐标函数读取变量；不需要把“无限次独立抽样”作为新的非数学原始对象。后文的大数定律将研究这些坐标的平均值，中心极限定理则研究其中心化和的分布。

% !TeX root = ../../Book1.tex
\section{Borel–Cantelli 引理与基本概率估计}
\label{b1:sec:3003}

积分提供平均大小，概率估计则把平均大小转化为偏差事件的上界。单次估计本身不涉及长期行为；当一列偏差概率可求和时，Borel--Cantelli 引理才把它们组织成几乎必然的最终结论。反向的无穷多次发生准则需要独立性，因而也说明了上一节的结构为何重要。

\ProseParagraphBreak

两条引理可对照 Axler 的《Measure, Integration \& Real Analysis》\cite[Borel--Cantelli Lemma 12.6, relative of Borel--Cantelli Lemma 12.10, pp.~383, 385]{Axler2020Measure}；第一条的一般测度空间版本另见 Fremlin 的《Measure Theory》第二卷\cite[Lemma 273A]{Fremlin2016Measure}。本节保留全部证明，并在最后用具体事件区分可求和、独立和概率趋零这三个条件。

\subsection{Markov 与 Chebyshev 估计}
\label{b1:subsec:300301}

\term[markov-budengshi]{Markov 不等式}[Markov's inequality]是\cref{b1:prop:integral-markov}的概率写法。它只需要积分的单调性，不要求变量独立，也不要求其分布有密度。

\begin{theorem}[Markov]\label{b1:thm:30-markov}
设 \((\Omega,\mathcal F,\mathbb P)\) 是概率空间，\(X\geq0\) 是实随机变量。则对每个 \(a>0\)，
\[
 \mathbb P(X\geq a)\leq\frac{\mathbb E X}{a}.
\]
更一般地，若 \(Z\) 是实值或复值可测函数，\(p>0\)，则对每个 \(a>0\)，
\[
 \mathbb P(|Z|\geq a)\leq\frac{\mathbb E|Z|^p}{a^p},
\]
其中期望可以为 \(+\infty\)。
\end{theorem}
\begin{proof}
对 \((\Omega,\mathcal F,\mathbb P)\) 上的非负可测函数 \(X\)，\cref{b1:prop:integral-markov}给出 \(a\cdot\mathbb P(X\geq a)\leq\mathbb E X\)。再将它用于 \(|Z|^p\) 与阈值 \(a^p\)，利用 \(\{|Z|\geq a\}=\{|Z|^p\geq a^p\}\)，即得第二式。
\end{proof}

例如，若 \(X=b\cdot\mathbf1_A\)，其中 \(b>0\)，则在阈值 \(a=b\) 处有 \(\mathbb P(X\geq b)=\mathbb E X/b=\mathbb P(A)\)。所以若只知道非负变量的期望，上述常数不能统一改小。若知道的是围绕均值的平方误差，就得到更适合平均值问题的估计。

\begin{corollary}[Chebyshev]\label{b1:cor:30-chebyshev}
设 \(X\in L^2(\mathbb P)\) 是概率空间 \((\Omega,\mathcal F,\mathbb P)\) 上的实随机变量。则对每个 \(a>0\)，
\[
 \mathbb P(|X-\mathbb E X|\geq a)
 \leq\frac{\operatorname{Var}(X)}{a^2}.
\]
若 \(\sigma_X=\sqrt{\operatorname{Var}(X)}>0\)，则对每个 \(t>0\)，还有
\[
 \mathbb P(|X-\mathbb E X|\geq t\sigma_X)\leq t^{-2}.
\]
\end{corollary}
\begin{proof}
在\cref{b1:thm:30-markov}中取 \(Z=X-\mathbb E X\)、\(p=2\)，并使用方差的定义。若 \(\sigma_X>0\)，再取 \(a=t\sigma_X>0\)，便得第二式。
\end{proof}

这个估计称为\term[chebyshev-budengshi]{Chebyshev 不等式}[Chebyshev's inequality]，可对照 Axler 的《Measure, Integration \& Real Analysis》\cite[Chebyshev's inequality 12.21, p.~389]{Axler2020Measure}。当 \(\sigma_X=0\) 时，\cref{b1:prop:30-variance-formula}已经给出 \(X=\mathbb E X\) 几乎必然，故任意正阈值的偏差概率为零。将阈值写成 \(t\sigma_X\) 时必须另作这个处理，因为零阈值下的事件 \(\{|X-\mathbb E X|\geq0\}\) 是整个空间。

\subsection{第一 Borel–Cantelli 引理}
\label{b1:subsec:300302}

设 \((A_n)_{n\geq1}\) 是概率空间 \((\Omega,\mathcal F,\mathbb P)\) 上的事件列。事件
\[
 \limsup_{n\to\infty}A_n
 =\bigcap_{m=1}^{\infty}\bigcup_{n=m}^{\infty}A_n
\]
由在无穷多个 \(A_n\) 中出现的样本点组成：无论从哪一项以后开始看，仍能遇到一次 \(A_n\)。其补集为
\[
 \bigcup_{m=1}^{\infty}\bigcap_{n=m}^{\infty}A_n^c,
\]
即只有有限多个 \(A_n\) 发生的事件。以下用“\(A_n\) 无穷多次发生”表示前一个事件。

\begin{lemma}[第一 Borel--Cantelli]\label{b1:lem:30-borel-cantelli-first}
设 \((\Omega,\mathcal F,\mathbb P)\) 是概率空间，\((A_n)_{n\geq1}\subseteq\mathcal F\) 是事件列。若
\[
 \sum_{n=1}^{\infty}\mathbb P(A_n)<\infty,
\]
则
\[
 \mathbb P\left(\limsup_{n\to\infty}A_n\right)=0.
\]
换言之，几乎必然只有有限多个 \(A_n\) 发生。
\end{lemma}
\begin{proof}
对每个 \(m\geq1\)，有
\[
 \limsup_nA_n\subseteq\bigcup_{n=m}^{\infty}A_n,
 \quad
 \mathbb P\left(\limsup_nA_n\right)
 \leq\sum_{n=m}^{\infty}\mathbb P(A_n).
\]
收敛级数的尾和趋零，故左端为零。
\end{proof}

这就是\term[borel-cantelli-yinli]{Borel--Cantelli 引理}[Borel--Cantelli lemma]的第一部分。证明没有使用独立性：可求和性本身已经使全部远期事件的并具有任意小的概率。它把一串数值估计提升为逐个样本的最终性质。

\begin{corollary}[可求和的误差事件]\label{b1:cor:30-summable-errors}
设 \(X_n,X\) 是概率空间 \((\Omega,\mathcal F,\mathbb P)\) 上的实随机变量。下列任一条件都推出 \(X_n\to X\) 几乎必然：
\begin{enumerate}
 \item 存在确定的正数列 \(a_n\to0\)，满足 \(\sum_{n\geq1}\mathbb P(|X_n-X|>a_n)<\infty\)。
 \item 对每个整数 \(k\geq1\)，都有 \(\sum_{n\geq1}\mathbb P(|X_n-X|>1/k)<\infty\)。
\end{enumerate}
\end{corollary}
\begin{proof}
在第一种情形，把第一引理用于 \(A_n=\{|X_n-X|>a_n\}\)。除去一个零测集后，存在依赖于样本点的 \(N\)，使 \(n\geq N\) 时 \(|X_n-X|\leq a_n\)。因为 \(a_n\to0\)，便得到逐点收敛。

在第二种情形，对每个 \(k\) 分别应用第一引理，得到零测集 \(N_k\)，使 \(\omega\notin N_k\) 时最终有 \(|X_n(\omega)-X(\omega)|\leq1/k\)。并集 \(N=\bigcup_{k\geq1}N_k\) 仍为零测集。在其补集上，给定 \(\varepsilon>0\)，先选 \(k\) 使 \(1/k<\varepsilon\)，再选上述最终指标，即得 \(X_n(\omega)\to X(\omega)\)。
\end{proof}

例如，若 \(X_n,X\) 在同一概率空间上，且对某个 \(p>0\) 有
\[
 \sum_{n=1}^{\infty}\mathbb E|X_n-X|^p<\infty,
\]
则由 Markov 不等式，对每个 \(k\geq1\) 有
\[
 \sum_{n=1}^{\infty}\mathbb P(|X_n-X|>1/k)
 \leq k^p\sum_{n=1}^{\infty}\mathbb E|X_n-X|^p<\infty.
\]
因此误差不仅依概率趋零，还几乎必然趋零。这里只要求误差矩的级数收敛，不要求变量间独立。

\subsection{独立性与第二 Borel–Cantelli 引理}
\label{b1:subsec:300303}

若 \(\sum_n\mathbb P(A_n)=\infty\)，尾并估计不再提供信息。独立性使我们能够改算它的补集：一长串事件全不发生的概率等于一串补概率的乘积，而发散的概率和将使该乘积趋零。

\begin{lemma}[第二 Borel--Cantelli]\label{b1:lem:30-borel-cantelli-second}
设 \((\Omega,\mathcal F,\mathbb P)\) 是概率空间，\((A_n)_{n\geq1}\subseteq\mathcal F\) 是相互独立的事件列。若
\[
 \sum_{n=1}^{\infty}\mathbb P(A_n)=\infty,
\]
则
\[
 \mathbb P\left(\limsup_{n\to\infty}A_n\right)=1.
\]
\end{lemma}
\begin{proof}
由\cref{b1:cor:30-events-generated-sigma}，补事件族也独立。对固定 \(m\geq1\) 与任意 \(N\geq m\)，有
\[
 \begin{aligned}
 \mathbb P\left(\bigcap_{n=m}^NA_n^c\right)
 &=\prod_{n=m}^N\bigl(1-\mathbb P(A_n)\bigr)\\
 &\leq\exp\left(-\sum_{n=m}^N\mathbb P(A_n)\right).
 \end{aligned}
\]
这里用到了 \(1-u\leq e^{-u}\)、\(0\leq u\leq1\)。当 \(N\to\infty\) 时，发散的尾和使右端趋零。左侧的事件递减，概率测度的上连续性给出
\[
 \mathbb P\left(\bigcap_{n=m}^{\infty}A_n^c\right)=0.
\]
对 \(m\) 取可数并，便有
\[
 \mathbb P\left(\bigcup_{m=1}^{\infty}\bigcap_{n=m}^{\infty}A_n^c\right)=0.
\]
括号内是只有有限多个 \(A_n\) 发生的事件，故所需结论成立。
\end{proof}

合并两条引理可知，对相互独立的事件列，无穷多次发生的概率由概率级数的敛散完全决定，分别为零或一。这里的零一结论只针对这个事件，并不是说独立试验中的每个事件都具有零一概率。

\ProseParagraphBreak

在独立同分布 Bernoulli 试验中，若 \(0<p<1\)，事件 \(\{X_n=1\}\) 的概率和与 \(\{X_n=0\}\) 的概率和都发散。分别应用第二引理，再把两个概率一事件相交，可知零与一几乎必然都出现无穷多次。这个结论只保证反复出现；出现频率是否趋于 \(p\)，是大数定律所回答的更细问题。

\subsection{几乎必然事件与反例}
\label{b1:subsec:300304}

先看独立性不可省略的例子。在 \(((0,1),\mathcal B((0,1)),m_1)\) 上取
\[
 A_n=(0,1/n).
\]
虽然 \(\sum_n\mathbb P(A_n)=\sum_n1/n=\infty\)，但这些事件递减，且其交为空，所以 \(\limsup_nA_n=\varnothing\)。每个样本点只属于有限多个 \(A_n\)。例如 \(\mathbb P(A_2\cap A_3)=1/3\)，而 \(\mathbb P(A_2)\cdot\mathbb P(A_3)=1/6\)，明确显示了独立性失败。这个例子也说明，第一引理的可求和条件是充分条件，并非无穷多次发生概率为零的必要条件。

\ProseParagraphBreak

无独立性时，发散概率和甚至不迫使无穷多次发生的概率取零或一。在 \([0,1]\) 的 Lebesgue 概率空间上，固定 \(A=[0,1/2]\)，令全部 \(A_n=A\)。此时概率和发散，但 \(\limsup_nA_n=A\) 的概率为 \(1/2\)。更一般地，重复一个概率为 \(q\in(0,1)\) 的事件，就会得到同样的 \(q\)。

\ProseParagraphBreak

最后，即使事件独立，单独知道 \(\mathbb P(A_n)\to0\) 也不能保证最终不再发生。由\cref{b1:cor:30-iid-construction}取独立 Bernoulli 变量 \(X_n\)，其中
\[
 \mathbb P(X_n=1)=\frac1{n+1},
 \quad \mathbb P(X_n=0)=1-\frac1{n+1}.
\]
对 \(A_n=\{X_n=1\}\)，有 \(\mathbb P(A_n)\to0\)，但概率和发散，所以第二引理给出 \(X_n=1\) 几乎必然出现无穷多次。另一方面，对每个有限 \(p\geq1\)，
\[
 \mathbb E|X_n|^p=\frac1{n+1}\longrightarrow0.
\]
故 \(X_n\) 依概率并在这些 \(L^p\) 中趋零，却不几乎必然趋零。依概率收敛允许小概率误差在无限多个时刻反复出现；可求和估计才把这些误差压缩到几乎必然只有有限次。下一节的大数定律将把这种差别落实到平均值序列上。

% !TeX root = ../../Book1.tex
\section{弱大数定律}
\label{b1:sec:3004}

把同一个随机试验重复许多次，样本平均应当接近期望。这里有两个不同的问题：
固定误差以后，偏离均值的概率是否趋于零；以及能否在同一个概率为一的集合上，
使整条样本平均序列收敛。本节先处理前一个问题。
有限二阶矩使平均误差可以直接用方差计算；只有一次矩时，则先截去尾部，
再分别控制有界部分和被舍去部分。

\ProseParagraphBreak

有限方差的论证可对照 Sheldon Axler 的《Measure, Integration \& Real Analysis》%
\cite[12.37--12.38, pp.~396--397]{Axler2020Measure}。
一次矩下的截断与平均收敛可与 David H. Fremlin 的《Measure Theory》%
\cite[273N(a)]{Fremlin2016Measure}比较；下面用固定截断和二阶矩估计直接证明所需结论。
% 实读：Axler2020Measure 印刷 pp.396--397，12.37--12.38；
% Fremlin2016Measure 作者公开 chap27 摘本 273N(a)。Bass2024Real 的
% Section 21.3（印刷 pp.276--278）仅作教学背景，不作为正式引用依据。

\subsection{独立和的期望与方差}
\label{b1:subsec:300401}

期望的有限可加性不需要独立性。方差还包含不同变量之间的交叉项，
独立性在这里的作用是使这些项在减去均值以后消失。

\begin{proposition}[期望与方差的有限加法]
\label{b1:prop:30-variance-sum}
设 \((\Omega,\mathcal F,\mathbb P)\) 为概率空间，\(N\geq1\)，
\(X_1,\ldots,X_N\) 为其上两两独立的实值随机变量，且
\(\mathbb E(X_j^2)<\infty\) 对 \(1\leq j\leq N\) 成立。
对实数 \(a_1,\ldots,a_N\)，有
\[
 \mathbb E\left(\sum_{j=1}^N a_jX_j\right)
 =\sum_{j=1}^N a_j\cdot\mathbb E X_j,
 \quad
 \operatorname{Var}\left(\sum_{j=1}^N a_jX_j\right)
 =\sum_{j=1}^N a_j^2\cdot\operatorname{Var}(X_j).
\]
第一项等式即使没有两两独立的假设也成立。
\end{proposition}
\begin{proof}
概率空间的总质量为 \(1\)，故 Cauchy--Schwarz 不等式给出
\(\mathbb E|X_j|\leq\bigl(\mathbb E(X_j^2)\bigr)^{1/2}<\infty\)。
期望的线性性随后给出第一项等式。
令 \(Z_j=X_j-\mathbb E X_j\)。这些变量仍两两独立，且均值为零。
对 \(i\ne j\)，由\cref{b1:prop:30-independent-expectations}，有
\[
 \mathbb E(Z_iZ_j)=(\mathbb E Z_i)\cdot(\mathbb E Z_j)=0.
\]
所用乘积可积；也可直接由
\[
 \mathbb E|Z_iZ_j|
 \leq\bigl(\mathbb E(Z_i^2)\bigr)^{1/2}
       \cdot\bigl(\mathbb E(Z_j^2)\bigr)^{1/2}<\infty
\]
核对。展开有限和的平方，得到
\[
 \begin{aligned}
 \operatorname{Var}\left(\sum_{j=1}^N a_jX_j\right)
 &=\mathbb E\left(\sum_{j=1}^N a_jZ_j\right)^2\\
 &=\sum_{j=1}^N a_j^2\cdot\mathbb E(Z_j^2)
   +2\sum_{1\leq i<j\leq N}a_ia_j\cdot\mathbb E(Z_iZ_j)\\
 &=\sum_{j=1}^N a_j^2\cdot\operatorname{Var}(X_j).
 \end{aligned}
\]
\end{proof}

取 \(a_j=1/N\)，可见平均的方差不是各项方差的平均，而是它们之和除以 \(N^2\)。
同分布再使所有单项方差相同，于是样本平均的方差缩小为单项方差的 \(1/N\)。
后面的极限定理首先来自这个计算。

\subsection{有限方差情形的弱大数律}
\label{b1:subsec:300402}

样本平均依概率收敛到共同均值的结论称为
\term[ruo-dashu-dinglv]{弱大数定律}[weak law of large numbers]。
先在二阶矩有限的条件下同时得到均方收敛和偏差概率的上界。

\begin{theorem}[有限方差的弱大数定律]
\label{b1:thm:30-finite-variance-weak-law}
设 \((\Omega,\mathcal F,\mathbb P)\) 为概率空间，\((X_n)_{n\geq1}\)
为其上的独立同分布实值随机变量列，且 \(\mathbb E(X_1^2)<\infty\)。
令 \(m=\mathbb E X_1\)、\(\sigma^2=\operatorname{Var}(X_1)\)、
\(S_n=\sum_{k=1}^nX_k\)。则
\[
 \mathbb E\left|\frac{S_n}{n}-m\right|^2=\frac{\sigma^2}{n},
 \quad
 \mathbb P\left(\left|\frac{S_n}{n}-m\right|\geq\varepsilon\right)
 \leq\frac{\sigma^2}{n\varepsilon^2}
 \quad(n\geq1,\ \varepsilon>0).
\]
因此 \(S_n/n\) 在 \(L^2(\mathbb P)\) 中且依概率收敛于 \(m\)。
\end{theorem}
\begin{proof}
同分布保证 \(\mathbb E X_k=m\)、\(\operatorname{Var}(X_k)=\sigma^2\)。
由\cref{b1:prop:30-variance-sum}，有
\[
 \mathbb E\left(\frac{S_n}{n}\right)=m,
 \quad
 \mathbb E\left|\frac{S_n}{n}-m\right|^2
 =\operatorname{Var}\left(\frac{S_n}{n}\right)
 =\frac1{n^2}\sum_{k=1}^n\operatorname{Var}(X_k)
 =\frac{\sigma^2}{n}.
\]
对实值平方可积变量 \(S_n/n\) 应用\cref{b1:cor:30-chebyshev}，
得到所述概率上界。即使 \(\sigma=0\)，这个计算也仍然成立。
\end{proof}

同分布不是上述方差方法的必要条件。只要总方差相对于 \(n^2\) 足够小，
平均和减去平均期望以后的误差也会消失。

\begin{corollary}[非同分布变量的方差条件]
\label{b1:cor:30-variance-weak-law}
设 \((\Omega,\mathcal F,\mathbb P)\) 为概率空间，\((X_n)_{n\geq1}\)
为其上两两独立的实值随机变量列，且每个 \(X_n\in L^2(\mathbb P)\)。
若
\[
 \lim_{n\to\infty}\frac1{n^2}
       \sum_{k=1}^n\operatorname{Var}(X_k)=0,
\]
则
\[
 \frac1n\sum_{k=1}^n(X_k-\mathbb E X_k)\longrightarrow0
 \quad\text{于 }L^2(\mathbb P)\text{ 且依概率成立}.
\]
若另有 \(n^{-1}\sum_{k=1}^n\mathbb E X_k\to m\in\mathbb R\)，
则 \(n^{-1}\sum_{k=1}^nX_k\) 在这两种意义下收敛于 \(m\)。
\end{corollary}
\begin{proof}
\cref{b1:prop:30-variance-sum}给出
\[
 \mathbb E\left|\frac1n\sum_{k=1}^n(X_k-\mathbb E X_k)\right|^2
 =\frac1{n^2}\sum_{k=1}^n\operatorname{Var}(X_k)\longrightarrow0.
\]
再用 Chebyshev 不等式得到依概率收敛。
若平均期望趋于 \(m\)，在左端随机变量上加上确定性的误差
\(n^{-1}\sum_{k=1}^n\mathbb E X_k-m\)，并用 \(L^2\) 三角不等式，即得最后的结论。
\end{proof}

例如，一致有界的单项方差满足这一条件。对独立的 Bernoulli 变量，若
\(\mathbb P(X_k=1)=p\)、\(\mathbb P(X_k=0)=1-p\)，其中 \(0\leq p\leq1\)，
则 \(S_n\) 是前 \(n\) 次试验中事件发生的次数，且
\[
 \mathbb P\left(\left|\frac{S_n}{n}-p\right|\geq\varepsilon\right)
 \leq\frac{p(1-p)}{n\varepsilon^2}
 \leq\frac1{4n\varepsilon^2}.
\]
这一上界控制任意固定样本量的偏差概率，但对 \(n\) 求和一般发散。
因此不能只把它代入第一 Borel--Cantelli 引理，就宣称整条样本平均序列几乎必然收敛。
下一节将另外控制一段部分和中的最大偏差。

\subsection{可积情形的截断与推广}
\label{b1:subsec:300403}

有限方差使计算简便，却不是弱大数定律所必需的。设一次矩有限，
先在一个固定高度截断所有变量。截断后的变量有有限方差，
而原变量与截断变量的平均差，可以直接用它们的一次矩尾部控制。
截断高度先固定，样本量趋于无穷以后，再让该高度增大。

\begin{theorem}[可积变量的弱大数定律]
\label{b1:thm:30-integrable-weak-law}
设 \((\Omega,\mathcal F,\mathbb P)\) 为概率空间，\((X_n)_{n\geq1}\)
为其上的独立同分布实值随机变量列，且 \(\mathbb E|X_1|<\infty\)。
令 \(m=\mathbb E X_1\)、\(S_n=\sum_{k=1}^nX_k\)。则
\[
 \lim_{n\to\infty}\mathbb E\left|\frac{S_n}{n}-m\right|=0,
 \quad
 \lim_{n\to\infty}
 \mathbb P\left(\left|\frac{S_n}{n}-m\right|\geq\varepsilon\right)=0
 \quad(\varepsilon>0).
\]
即样本平均在 \(L^1(\mathbb P)\) 中且依概率收敛于共同均值。
\end{theorem}
\begin{proof}
固定实数 \(R\geq1\)，令
\[
 Y_k^{(R)}=X_k\cdot\mathbf1_{\{|X_k|\leq R\}},
 \quad
 m_R=\mathbb E Y_1^{(R)},
 \quad
 \delta_R=\mathbb E\bigl(|X_1|\cdot\mathbf1_{\{|X_1|>R\}}\bigr).
\]
同一个 Borel 函数作用在各个 \(X_k\) 上，所得 \(Y_k^{(R)}\) 仍独立同分布。
它们满足 \(|Y_k^{(R)}|\leq R\)，所以
\(\operatorname{Var}(Y_1^{(R)})\leq\mathbb E((Y_1^{(R)})^2)\leq R^2\)。
由有限方差的弱大数定律及 Cauchy--Schwarz 不等式，
\[
 \mathbb E\left|\frac1n\sum_{k=1}^nY_k^{(R)}-m_R\right|
 \leq\left(
  \mathbb E\left|\frac1n\sum_{k=1}^nY_k^{(R)}-m_R\right|^2
       \right)^{1/2}
 \leq\frac R{\sqrt n}.
\]
另一方面，同分布给出
\[
 \mathbb E|X_k-Y_k^{(R)}|=\delta_R,
 \quad |m-m_R|\leq\delta_R.
\]
把样本平均的误差分成三项，再取期望，得到
\[
 \begin{aligned}
 \mathbb E\left|\frac{S_n}{n}-m\right|
 &\leq\mathbb E\left|\frac1n\sum_{k=1}^nY_k^{(R)}-m_R\right|
       +\frac1n\sum_{k=1}^n\mathbb E|X_k-Y_k^{(R)}|+|m_R-m|\\
 &\leq\frac R{\sqrt n}+2\delta_R.
 \end{aligned}
\]
先令 \(n\to\infty\)，得到左端的上极限至多为 \(2\delta_R\)。
因为 \(|X_1|\) 可积，且尾部被积函数由 \(|X_1|\) 控制并逐点趋于零，
控制收敛定理说明 \(\delta_R\to0\)。再令 \(R\to\infty\)，便得 \(L^1\) 收敛。
最后，由 Markov 不等式，
\[
 \mathbb P\left(\left|\frac{S_n}{n}-m\right|\geq\varepsilon\right)
 \leq\frac1\varepsilon\cdot
       \mathbb E\left|\frac{S_n}{n}-m\right|\longrightarrow0.
\]
\end{proof}

这个证明中，固定截断使二阶矩估计里的 \(R\) 在 \(n\to\infty\) 时保持不变。
若立即取 \(R=n\)，仅有一次矩时，上面的 \(R/\sqrt n\) 就不再趋于零，
需要另外的论证才能使用随样本量变化的截断。

\ProseParagraphBreak

可积而不平方可积的分布确实存在。在 \((0,1)\) 上取 Lebesgue 概率测度，并令
\(X(u)=u^{-2/3}\)，则
\[
 \mathbb E X=\int_0^1u^{-2/3}\,\mathrm du=3,
 \quad
 \mathbb E(X^2)=\int_0^1u^{-4/3}\,\mathrm du=\infty.
\]
取具有这个分布的独立随机变量列，其样本平均仍在 \(L^1\) 中且依概率趋于 \(3\)。
方差公式不能用于这列变量，而固定截断仍能控制它们的平均误差。

% !TeX root = ../../Book1.tex
\section{可积独立同分布变量的强大数定律}
\label{b1:sec:3005}

弱大数定律逐个控制样本量 \(n\) 对应的偏差事件。要证明几乎必然收敛，
则须在同一个概率为一的集合上控制所有充分大的样本量。
本节先建立独立部分和的极大估计，由此证明方差可和的独立零均值级数几乎必然收敛，
再以求和变换把级数收敛转为平均收敛。一次矩有限的变量经过随指标增长的截断，
恰好满足这条级数判据。

\ProseParagraphBreak

独立级数、归一化求和与可积独立同分布大数定律，可对照 Fremlin 的
《Measure Theory》\cite[273B(b), 273C--273D, 273I]{Fremlin2016Measure}。
下面用部分和第一次越过阈值的事件直接证明所需的极大不等式，
再逐项核对截断后的方差、尾概率和均值。
% 实读：Fremlin2016Measure 作者公开 chap27 摘本 273B(b)、273C、273D、
% 273I（章摘本物理 pp.23--29），以及 273N(a)。正文改从 n=1 编号，
% 以 Kolmogorov 极大不等式自证独立级数收敛，不借该书 273B(a) 的 Etemadi 路线。
% Bass2024Real 的 Section 21.4（印刷 pp.278--282）仅作教学背景；
% 本节的加权截断矩估计直接对共同分布应用非负级数换序。

\subsection{独立部分和的极大估计}
\label{b1:subsec:300501}

对某一个 \(S_n\) 使用 Chebyshev 不等式，不能同时控制前面的每个部分和。
下面的\term[kolmogorov-jida-budengshi]{Kolmogorov 极大不等式}[Kolmogorov's maximal inequality]
用最后一个和的总方差，控制整段部分和是否越过固定阈值。

\begin{theorem}[Kolmogorov 极大不等式]
\label{b1:thm:30-kolmogorov-maximal}
设 \((\Omega,\mathcal F,\mathbb P)\) 为概率空间，\(N\geq1\)，
\(Z_1,\ldots,Z_N\) 为其上独立的实值随机变量，满足
\(\mathbb E Z_k=0\)、\(\mathbb E(Z_k^2)<\infty\)。令
\(S_j=\sum_{k=1}^jZ_k\)。则对每个 \(\lambda>0\)，有
\[
 \mathbb P\left(\max_{1\leq j\leq N}|S_j|\geq\lambda\right)
 \leq\frac1{\lambda^2}\sum_{k=1}^N\mathbb E(Z_k^2).
\]
\end{theorem}
\begin{proof}
对 \(1\leq j\leq N\)，令
\[
 A_j=\{\,|S_i|<\lambda\ (1\leq i<j),\quad |S_j|\geq\lambda\,\}.
\]
这些事件两两不交，且它们的并就是最大值至少为 \(\lambda\) 的事件。
事件 \(A_j\) 以及变量 \(S_j\) 只依赖 \(Z_1,\ldots,Z_j\)，
而 \(S_N-S_j\) 只依赖余下的变量。
\cref{b1:prop:30-independent-blocks}保证两个指标块生成的
\(\sigma\)-代数独立。因此
\(S_j\cdot\mathbf1_{A_j}\) 与 \(S_N-S_j\) 独立。
两者均平方可积，后一变量的均值为零，由独立变量的期望乘积公式，
\[
 \mathbb E\bigl(S_j\cdot\mathbf1_{A_j}\cdot(S_N-S_j)\bigr)=0.
\]
当 \(j=N\) 时，后一变量恒为零，同一等式也成立。
在 \(A_j\) 上展开 \(S_N=S_j+(S_N-S_j)\) 的平方，得到
\[
 \begin{aligned}
 \mathbb E\bigl(S_N^2\cdot\mathbf1_{A_j}\bigr)
 &=\mathbb E\bigl(S_j^2\cdot\mathbf1_{A_j}\bigr)
   +\mathbb E\bigl((S_N-S_j)^2\cdot\mathbf1_{A_j}\bigr)\\
 &\geq\lambda^2\cdot\mathbb P(A_j).
 \end{aligned}
\]
对不交的 \(A_j\) 求和，并用\cref{b1:prop:30-variance-sum}，有
\[
 \begin{aligned}
 \lambda^2\cdot
 \mathbb P\left(\max_{1\leq j\leq N}|S_j|\geq\lambda\right)
 &\leq\sum_{j=1}^N\mathbb E\bigl(S_N^2\cdot\mathbf1_{A_j}\bigr)\\
 &\leq\mathbb E(S_N^2)
 =\sum_{k=1}^N\mathbb E(Z_k^2).
 \end{aligned}
\]
除以 \(\lambda^2\) 即得结论。
\end{proof}

这里按第一次越过阈值的时刻划分事件，使每个样本至多计入一次。
对各个 \(S_j\) 分别应用 Chebyshev 不等式再相加，会反复计入前面变量的方差，
不能给出上面的同一总方差界。

\subsection{截断与尾概率控制}
\label{b1:subsec:300502}

强大数定律将使用截断 \(Y_n=X_n\cdot\mathbf1_{\{|X_n|\leq n\}}\)。
截断高度随 \(n\) 增长，有两个需要分别证明的后果：原变量最终不再被改变；
截断后方差除以 \(n^2\) 得到的级数可和。二者都由同一个有限一次矩控制。

\begin{lemma}[一次矩控制的截断估计]
\label{b1:lem:30-truncation-sums}
设 \((\Omega,\mathcal F,\mathbb P)\) 为概率空间，\(X\) 为其上的实值随机变量，
且 \(\mathbb E|X|<\infty\)。则
\[
 \sum_{n=1}^\infty\mathbb P(|X|>n)\leq\mathbb E|X|,
 \quad
 \sum_{n=1}^\infty\frac1{n^2}\cdot
       \mathbb E\bigl(|X|^2\cdot\mathbf1_{\{|X|\leq n\}}\bigr)
 \leq2\mathbb E|X|.
\]
\end{lemma}
\begin{proof}
对每个实数 \(u\geq0\)，小于 \(u\) 的正整数的个数不超过 \(u\)，故
\[
 \sum_{n=1}^\infty\mathbf1_{\{u>n\}}\leq u.
\]
取 \(u=|X|\) 并积分，非负级数的 Tonelli 换序给出第一项估计。

再令 \(m\) 为大于等于 \(u\) 的最小正整数。由递减函数的积分比较，
\[
 \sum_{n=m}^\infty\frac1{n^2}
 \leq\frac1{m^2}+\int_m^\infty\frac1{t^2}\,\mathrm dt
 \leq\frac2m.
\]
因为 \(m\geq u\)，对 \(u\geq0\) 有
\[
 \sum_{n=1}^\infty\frac{u^2}{n^2}\cdot\mathbf1_{\{u\leq n\}}
 =u^2\cdot\sum_{n=m}^\infty\frac1{n^2}
 \leq\frac{2u^2}{m}\leq2u.
\]
再次取 \(u=|X|\) 并积分。所有项均非负，故 Tonelli 换序不需要事先假设
\(X\) 平方可积，得到第二项估计。
\end{proof}

\begin{proposition}[随指标增长的截断]
\label{b1:prop:30-moving-truncation}
设 \((\Omega,\mathcal F,\mathbb P)\) 为概率空间，\((X_n)_{n\geq1}\)
为其上的独立同分布实值随机变量列，且 \(\mathbb E|X_1|<\infty\)。
令 \(Y_n=X_n\cdot\mathbf1_{\{|X_n|\leq n\}}\)。则 \((Y_n)\) 独立，
每个 \(Y_n\) 平方可积，且
\[
 \sum_{n=1}^\infty\mathbb P(X_n\ne Y_n)<\infty,
 \quad
 \sum_{n=1}^\infty\frac{\operatorname{Var}(Y_n)}{n^2}
 \leq2\mathbb E|X_1|<\infty.
\]
此外，\(\mathbb E Y_n\to\mathbb E X_1\)。
\end{proposition}
\begin{proof}
各个 \(Y_n\) 分别是对应 \(X_n\) 的 Borel 函数，故保持独立性。
由 \(|Y_n|\leq n\)，它们均平方可积。
同分布和前一引理给出
\[
 \sum_{n=1}^\infty\mathbb P(X_n\ne Y_n)
 =\sum_{n=1}^\infty\mathbb P(|X_1|>n)<\infty.
\]
又因为 \(\operatorname{Var}(Y_n)\leq\mathbb E(Y_n^2)\)，且同分布保证
\[
 \mathbb E(Y_n^2)
 =\mathbb E\bigl(|X_1|^2\cdot\mathbf1_{\{|X_1|\leq n\}}\bigr),
\]
第二个级数也由前一引理控制。最后，
\[
 \mathbb E Y_n
 =\mathbb E\bigl(X_1\cdot\mathbf1_{\{|X_1|\leq n\}}\bigr)
 \longrightarrow\mathbb E X_1,
\]
这里的支配函数为可积的 \(|X_1|\)。
\end{proof}

这些截断变量一般不再同分布，因而不能对它们直接套用独立同分布版本的大数定律。
下面的收敛工具只要求独立和方差可和，正好适用于这一列变量。

\subsection{归一化和的收敛工具}
\label{b1:subsec:300503}

先把有限部分和的极大估计用于级数尾部。方差尾和趋于零本身只给概率估计；
选择一列足够小的尾和，再应用第一 Borel--Cantelli 引理，才能得到逐样本的 Cauchy 性。

\begin{theorem}[独立零均值级数的收敛]
\label{b1:thm:30-independent-series}
设 \((\Omega,\mathcal F,\mathbb P)\) 为概率空间，\((Z_n)_{n\geq1}\)
为其上独立的实值随机变量列，且 \(\mathbb E Z_n=0\)、
\(\mathbb E(Z_n^2)<\infty\) 对每个 \(n\) 成立。若
\[
 \sum_{n=1}^\infty\mathbb E(Z_n^2)<\infty,
\]
则级数 \(\sum_{n=1}^\infty Z_n\) 几乎必然收敛到有限实值随机变量 \(S\)，
并在 \(L^2(\mathbb P)\) 中收敛。取 \(S_n=\sum_{k=1}^nZ_k\)，还有
\[
 \mathbb E S=0,
 \quad
 \mathbb E|S-S_n|^2=\sum_{k=n+1}^\infty\mathbb E(Z_k^2).
\]
\end{theorem}
\begin{proof}
令 \(v_m=\sum_{k=m}^\infty\mathbb E(Z_k^2)\)。对有限指标段 \(m,\ldots,N\)
应用 Kolmogorov 极大不等式，得到
\[
 \mathbb P\left(\max_{m\leq n\leq N}
       \left|\sum_{k=m}^nZ_k\right|>\varepsilon\right)
 \leq\frac{v_m}{\varepsilon^2}.
\]
当 \(N\to\infty\) 时，左侧事件递增，故
\[
 \mathbb P\left(\sup_{n\geq m}
       \left|\sum_{k=m}^nZ_k\right|>\varepsilon\right)
 \leq\frac{v_m}{\varepsilon^2}.
\]
这里的上确界只取可数多个可测函数，因此事件可测。
由于 \(v_m\to0\)，可选严格递增的正整数 \(m_j\)，使
\(v_{m_j}\leq2^{-3j}\)。令
\[
 H_j=\left\{\sup_{n\geq m_j}
          \left|\sum_{k=m_j}^nZ_k\right|>2^{-j}\right\}.
\]
由上一估计，\(\mathbb P(H_j)\leq2^{-j}\)。
第一 Borel--Cantelli 引理\cref{b1:lem:30-borel-cantelli-first}说明，
几乎每个样本只属于有限多个 \(H_j\)。固定这样的样本 \(\omega\)。
对充分大的 \(j\)，以及任意 \(n\geq p\geq m_j\)，有
\[
 \left|\sum_{k=p}^nZ_k(\omega)\right|
 \leq\left|\sum_{k=m_j}^nZ_k(\omega)\right|
     +\left|\sum_{k=m_j}^{p-1}Z_k(\omega)\right|
 \leq2^{1-j},
\]
其中空和取零。故 \((S_n(\omega))\) 是实数 Cauchy 列，收敛到有限值。
在余下的可测零测集上把极限置为零，得到实值随机变量 \(S\)。

还需核对 \(L^2\) 结论。固定 \(n\geq1\)，由几乎必然收敛和 Fatou 引理，
\[
 \mathbb E|S-S_n|^2
 \leq\liminf_{N\to\infty}\mathbb E|S_N-S_n|^2
 =\sum_{k=n+1}^\infty\mathbb E(Z_k^2).
\]
右端有限，故 \(S\in L^2(\mathbb P)\)；它随 \(n\to\infty\) 趋于零，
说明 \(S_n\to S\) 于 \(L^2\)。因此对固定 \(n\)，
\(S_N-S_n\to S-S_n\) 于 \(L^2\)，其范数也收敛。
在有限和的方差等式中令 \(N\to\infty\)，得到所述精确尾和公式。
又由 Cauchy--Schwarz 不等式，\(L^2\) 收敛蕴含 \(L^1\) 收敛，
所以 \(\mathbb E S=\lim_n\mathbb E S_n=0\)。
\end{proof}

接下来需要一个确定性的求和事实。
\term[kronecker-yinli]{Kronecker 引理}[Kronecker's lemma]
把除以 \(b_n\) 后形成的级数，转为原部分和除以 \(b_n\) 的极限。

\begin{lemma}[Kronecker]
\label{b1:lem:30-kronecker}
设 \(\mathbb K=\mathbb R\) 或 \(\mathbb C\)，\((a_n)_{n\geq1}\subset\mathbb K\)，
\((b_n)_{n\geq1}\) 为非递减的正实数列，且 \(b_n\to\infty\)。
若级数 \(\sum_{n=1}^\infty a_n/b_n\) 在 \(\mathbb K\) 中收敛，则
\[
 \lim_{n\to\infty}\frac1{b_n}\sum_{k=1}^na_k=0.
\]
\end{lemma}
\begin{proof}
令 \(s_n=\sum_{k=1}^na_k/b_k\)、\(s_0=0\)，并记 \(s=\lim_ns_n\)。
由 \(a_k=b_k(s_k-s_{k-1})\)，有限求和变换给出
\[
 \sum_{k=1}^na_k
 =b_ns_n-\sum_{k=1}^{n-1}(b_{k+1}-b_k)s_k.
\]
因为 \(\sum_{k=1}^{n-1}(b_{k+1}-b_k)=b_n-b_1\)，所以
\[
 \frac1{b_n}\sum_{k=1}^na_k
 =(s_n-s)+\frac{b_1}{b_n}s
   -\frac1{b_n}\sum_{k=1}^{n-1}(b_{k+1}-b_k)(s_k-s).
\]
前两项趋于零。对任意 \(\varepsilon>0\)，取 \(M\) 使
\(|s_k-s|\leq\varepsilon\) 对 \(k\geq M\) 成立。
末项中 \(k<M\) 的有限和除以 \(b_n\) 后趋于零；其余项满足
\[
 \frac1{b_n}\sum_{k=M}^{n-1}(b_{k+1}-b_k)\cdot|s_k-s|
 \leq\varepsilon\cdot\frac{b_n-b_M}{b_n}\leq\varepsilon.
\]
这里用了 \(b_{k+1}-b_k\geq0\)。再令 \(\varepsilon\downarrow0\)，即得结论。
\end{proof}

\begin{corollary}[方差条件下的归一化强大数定律]
\label{b1:cor:30-variance-strong-law}
设 \((\Omega,\mathcal F,\mathbb P)\) 为概率空间，\((X_n)_{n\geq1}\)
为其上独立的实值随机变量列，每个 \(X_n\in L^2(\mathbb P)\)。
设 \((b_n)_{n\geq1}\) 为非递减的正实数列，\(b_n\to\infty\)，且
\[
 \sum_{n=1}^\infty\frac{\operatorname{Var}(X_n)}{b_n^2}<\infty.
\]
则
\[
 \frac1{b_n}\sum_{k=1}^n(X_k-\mathbb E X_k)
 \longrightarrow0\quad\text{几乎必然成立}.
\]
特别地，若 \(\sum_{n=1}^\infty\operatorname{Var}(X_n)/n^2<\infty\)，
则可取 \(b_n=n\)。
\end{corollary}
\begin{proof}
令 \(Z_n=(X_n-\mathbb E X_n)/b_n\)。这些变量独立、均值为零，且
\[
 \sum_{n=1}^\infty\mathbb E(Z_n^2)
 =\sum_{n=1}^\infty\frac{\operatorname{Var}(X_n)}{b_n^2}<\infty.
\]
由\cref{b1:thm:30-independent-series}，级数 \(\sum_nZ_n\) 在一个概率为一的集合上收敛。
在这个集合的每一点，对确定性数列 \(a_n=X_n(\omega)-\mathbb E X_n\)
应用 Kronecker 引理，便得到所述极限。
\end{proof}

这个结论允许各项的分布不同。例如，独立变量的方差若一致有界，则
\[
 \sum_{n=1}^{\infty}\frac{\operatorname{Var}(X_n)}{n^2}<\infty.
\]
因此，平均和减去平均期望以后几乎必然趋于零。

\subsection{可积情形的强大数律与边界}
\label{b1:subsec:300504}

最后还要把逐项均值的收敛传到它们的平均。
对于数列 \((c_n)\)，\(n^{-1}\sum_{k=1}^nc_k\) 称为它的
\term[cesaro-pingjun]{Cesàro 平均}[Cesàro mean]。

\begin{lemma}[Cesàro 平均]
\label{b1:lem:30-cesaro}
设 \(\mathbb K=\mathbb R\) 或 \(\mathbb C\)，\((c_n)_{n\geq1}\subset\mathbb K\)，
且 \(c_n\to c\in\mathbb K\)。则
\[
 \frac1n\sum_{k=1}^nc_k\longrightarrow c.
\]
\end{lemma}
\begin{proof}
给定 \(\varepsilon>0\)，取 \(M\) 使 \(|c_k-c|\leq\varepsilon\) 对 \(k\geq M\) 成立。
对 \(n\geq M\)，有
\[
 \left|\frac1n\sum_{k=1}^nc_k-c\right|
 \leq\frac1n\sum_{k=1}^{M-1}|c_k-c|
       +\frac{n-M+1}{n}\cdot\varepsilon.
\]
第一项是固定有限和除以 \(n\)，趋于零；第二项至多为 \(\varepsilon\)。
令 \(\varepsilon\downarrow0\) 即得结论。
\end{proof}

样本平均几乎必然收敛到共同均值的结论称为
\term[qiang-dashu-dinglv]{强大数定律}[strong law of large numbers]。
对于独立同分布的实值变量，一个有限一次矩便足够。

\begin{theorem}[可积独立同分布变量的强大数定律]
\label{b1:thm:30-iid-strong-law}
设 \((\Omega,\mathcal F,\mathbb P)\) 为概率空间，\((X_n)_{n\geq1}\)
为其上的独立同分布实值随机变量列，且 \(\mathbb E|X_1|<\infty\)。
令 \(m=\mathbb E X_1\)、\(S_n=\sum_{k=1}^nX_k\)。则
\[
 \mathbb P\left(\left\{\omega\in\Omega:
       \lim_{n\to\infty}\frac{S_n(\omega)}n=m\right\}\right)=1.
\]
\end{theorem}
\begin{proof}
取 \(Y_n=X_n\cdot\mathbf1_{\{|X_n|\leq n\}}\)。
\cref{b1:prop:30-moving-truncation}保证 \((Y_n)\) 独立，每个 \(Y_n\) 平方可积，且
\[
 \sum_{n=1}^\infty\frac{\operatorname{Var}(Y_n)}{n^2}<\infty.
\]
因此\cref{b1:cor:30-variance-strong-law}在 \(b_n=n\) 时给出
\[
 \frac1n\sum_{k=1}^n(Y_k-\mathbb E Y_k)\longrightarrow0
 \quad\text{几乎必然成立}.
\]
同一截断命题还给出 \(\sum_n\mathbb P(X_n\ne Y_n)<\infty\)。
由第一 Borel--Cantelli 引理，几乎每个 \(\omega\) 都有一个正整数 \(N(\omega)\)，
使 \(X_n(\omega)=Y_n(\omega)\) 对 \(n\geq N(\omega)\) 成立。
所以在这个概率为一的集合上，对 \(n\geq N(\omega)\) 有
\[
 \frac1n\sum_{k=1}^n(X_k(\omega)-Y_k(\omega))
 =\frac1n\sum_{k=1}^{N(\omega)-1}(X_k(\omega)-Y_k(\omega))
 \longrightarrow0.
\]
这是一个固定有限的实数和除以 \(n\)，不需要对未截断变量作二阶矩估计。

最后，\(\mathbb E Y_k\to m\) 已由截断命题中的控制收敛证明。
Cesàro 平均引理于是给出
\[
 \frac1n\sum_{k=1}^n\mathbb E Y_k\longrightarrow m.
\]
在前两个几乎必然结论的共同概率为一的集合上，分解
\[
 \frac{S_n}{n}-m
 =\frac1n\sum_{k=1}^n(X_k-Y_k)
  +\frac1n\sum_{k=1}^n(Y_k-\mathbb E Y_k)
  +\left(\frac1n\sum_{k=1}^n\mathbb E Y_k-m\right),
\]
三项分别趋于零，得到所述收敛。
\end{proof}

这与\cref{b1:thm:30-integrable-weak-law}相合：可积独立同分布样本的平均，
既几乎必然收敛，也在 \(L^1\) 中收敛。极大估计控制整条样本路径，固定截断则控制积分误差，两种方法分别给出上述两种收敛。

\ProseParagraphBreak

一次矩条件对于这里的几乎必然有限极限还是必要的。证明需要反过来观察单项 \(X_n/n\)，
并使用第二 Borel--Cantelli 引理。

\begin{proposition}[有限极限所需的一次矩]
\label{b1:prop:30-strong-law-integrability-necessary}
设 \((\Omega,\mathcal F,\mathbb P)\) 为概率空间，\((X_n)_{n\geq1}\)
为其上的独立同分布实值随机变量列，\(S_n=\sum_{k=1}^nX_k\)。
若存在实值随机变量 \(L\)，使 \(S_n/n\to L\) 几乎必然成立，
则 \(\mathbb E|X_1|<\infty\)，且 \(L=\mathbb E X_1\) 几乎必然成立。
\end{proposition}
\begin{proof}
对 \(n\geq2\)，有
\[
 \frac{X_n}{n}
 =\frac{S_n}{n}-\frac{n-1}{n}\cdot\frac{S_{n-1}}{n-1}
 \longrightarrow L-L=0\quad\text{几乎必然成立}.
\]
因此事件 \(A_n=\{|X_n|>n\}\) 几乎必然只发生有限次。
这些事件相互独立；若 \(\sum_n\mathbb P(A_n)=\infty\)，
第二 Borel--Cantelli 引理\cref{b1:lem:30-borel-cantelli-second}会给出它们几乎必然发生无穷多次，矛盾。
所以
\[
 \sum_{n=1}^\infty\mathbb P(|X_1|>n)
 =\sum_{n=1}^\infty\mathbb P(A_n)<\infty.
\]
对每个有限实数 \(u\geq0\)，有
\(u\leq1+\sum_{n=1}^\infty\mathbf1_{\{u>n\}}\)。
取 \(u=|X_1|\) 并以 Tonelli 定理换序，得到
\[
 \mathbb E|X_1|
 \leq1+\sum_{n=1}^\infty\mathbb P(|X_1|>n)<\infty.
\]
于是可应用已经证明的强大数定律；极限的唯一性给出
\(L=\mathbb E X_1\) 几乎必然成立。
\end{proof}

独立性也不能直接删除。在 \(\Omega=\{0,1\}\) 上令两个点的概率均为 \(1/2\)，
并令所有 \(X_n(\omega)=\omega\)。这些变量有界、同分布且期望均为 \(1/2\)，
但它们完全相同，故
\[
 \frac1n\sum_{k=1}^nX_k=X_1,
 \quad
 \mathbb P\left(\left|\frac1n\sum_{k=1}^nX_k-\frac12\right|>\frac14\right)=1.
\]
样本平均甚至不依概率收敛到均值。独立的重复观测与同一个观测的重复记录，
在平均极限中有不同的行为。

\ProseParagraphBreak

大数定律确定了平均的极限，却没有给出围绕均值波动的分布。
有限方差情形的 \(1/n\) 方差计算提示，平均误差的自然尺度为 \(1/\sqrt n\)。
后面的特征函数与中心极限定理将研究这一尺度上的分布极限。

% !TeX root = ../../Book1.tex
% 实读来源：Fremlin, Measure Theory, Vol. 2，作者在线第28章摘本，
% 285A、285F、285I、285J、285L--M；摘本 PDF 第67、69--73页。
% 条号作为正式定位，不把持续修订摘本的 PDF 页码当作2016版页码。
% Bass, Real Analysis for Graduate Students, v5，印刷第297--302、305--307页，
% 仅作教学背景。以下 Gaussian 平滑唯一性与极限存在性均在正文自证。
\section{特征函数与概率弱收敛}
\label{b1:sec:3006}

大数律比较随机变量与一个常数的距离；中心极限定理则要比较整个分布。
直接对所有有界连续函数计算积分，通常不便于处理独立和。
复指数函数的优势在于：它始终有界，并且把加法变成乘法。
本节先用这些特殊的测试函数表示分布，再说明它们为什么足以辨认分布及其弱极限。

\subsection{特征函数与 Gaussian 模型}
\label{b1:subsec:300601}

对实数 \(u\)，记 \(e^{iu}=\cos u+i\sin u\)。
这里的复值积分按实部、虚部分别定义；由于 \(|e^{iu}|=1\)，
对任意概率分布都可以积分，不需要假设一阶矩存在。

\begin{definition}\label{b1:def:30-characteristic-function}
设 \(\mu\) 是 \(\mathbb R\) 上的 Borel 概率测度。
函数
\[
 \varphi_\mu(t)=\int_{\mathbb R}e^{itx}\,\mathrm d\mu(x),
 \quad t\in\mathbb R,
\]
称为 \(\mu\) 的\term[tezheng-hanshu]{特征函数}[characteristic function]。
若 \(X\) 是概率空间 \((\Omega,\mathcal F,\mathbb P)\) 上的实随机变量，
其分布为 \(\mu_X\)，则规定
\[
 \varphi_X(t)=\varphi_{\mu_X}(t)=\mathbb E(e^{itX}).
\]
\end{definition}
\symidx{\(\varphi_\mu,\varphi_X\)：概率分布或随机变量的特征函数}

这一约定在指数中取正号，且不加归一化系数。
特征函数的定义和下面的基本运算可对照
Fremlin 的《Measure Theory》\cite[285A, 285F, 285I]{Fremlin2016Measure}。

\begin{proposition}\label{b1:prop:30-characteristic-basic}
设 \(\mu\) 是实轴上的 Borel 概率测度，则
\[
 \varphi_\mu(0)=1,\quad
 |\varphi_\mu(t)|\leq1,\quad
 \varphi_\mu(-t)=\overline{\varphi_\mu(t)}
 \quad(t\in\mathbb R),
\]
且 \(\varphi_\mu\) 在 \(\mathbb R\) 上一致连续。
若 \(X\) 是任意概率空间上的实随机变量，\(a,b\in\mathbb R\)，则
\[
 \varphi_{aX+b}(t)=e^{itb}\cdot\varphi_X(at).
\]
若 \(X_1,\ldots,X_r\) 是同一概率空间上的独立实随机变量，
其中整数 \(r\geq1\)，则
\begin{equation}\label{b1:eq:30-independent-characteristic}
 \varphi_{X_1+\cdots+X_r}(t)
 =\prod_{j=1}^r\varphi_{X_j}(t)
 \quad(t\in\mathbb R).
\end{equation}
\end{proposition}
\begin{proof}
前三式分别由总质量为一、积分三角不等式和共轭与积分交换得到。
对任意 \(t,h\in\mathbb R\)，有
\[
 |\varphi_\mu(t+h)-\varphi_\mu(t)|
 \leq\int_{\mathbb R}|e^{ihx}-1|\,\mathrm d\mu(x).
\]
当 \(h\to0\) 时，被积函数逐点趋于零，并由可积常数 \(2\) 控制。
控制收敛定理使右端趋于零；右端又与 \(t\) 无关，故得到一致连续性。

仿射变换公式来自
\(e^{it(aX+b)}=e^{itb}\cdot e^{i(at)X}\)。
对独立变量，\cref{b1:prop:30-independent-joint-law}给出联合分布
\(\mu_{X_1}\otimes\cdots\otimes\mu_{X_r}\)。
被积函数 \(\prod_{j=1}^r e^{itx_j}\) 的模恒为一，
因此在这个乘积概率空间上绝对可积。
Fubini 定理遂将其积分拆成各个特征函数的乘积，
得到\eqref{b1:eq:30-independent-characteristic}。
\end{proof}

单点分布 \(\delta_a\) 的特征函数是 \(e^{ita}\)；
参数 \(p\in[0,1]\) 的 Bernoulli 分布的特征函数是
\(1-p+p e^{it}\)。
接下来需要一个连续分布模型，它的特征函数仍具有便于取幂的形式。
Gaussian 密度具有这一性质，下面计算它的特征函数。

\begin{lemma}\label{b1:lem:30-gaussian-fourier}
对每个 \(\varepsilon>0\)，定义
\[
 g_\varepsilon(x)=\frac{1}{\sqrt{2\pi}\,\varepsilon}
       \exp\left(-\frac{x^2}{2\varepsilon^2}\right),
 \quad x\in\mathbb R.
\]
则 \(g_\varepsilon\) 是严格正的连续概率密度，并满足
\[
 \int_{\mathbb R}x\cdot g_\varepsilon(x)\,\mathrm dx=0,
 \quad
 \int_{\mathbb R}x^2\cdot g_\varepsilon(x)\,\mathrm dx=\varepsilon^2.
\]
此外，对所有 \(t,x\in\mathbb R\)，有
\begin{align}
 \int_{\mathbb R}e^{itx}\cdot g_\varepsilon(x)\,\mathrm dx
   &=e^{-\varepsilon^2t^2/2},
       \label{b1:eq:30-gaussian-characteristic}\\
 g_\varepsilon(x)
   &=\frac{1}{2\pi}\int_{\mathbb R}
       e^{-itx}\cdot e^{-\varepsilon^2t^2/2}\,\mathrm dt.
       \label{b1:eq:30-gaussian-inversion}
\end{align}
上述各积分均绝对收敛。
\end{lemma}
\begin{proof}
\secref{b1:sec:0705}已用 Tonelli 定理和极坐标换元证明
\(\int_{\mathbb R}e^{-u^2}\,\mathrm du=\sqrt\pi\)。
在此公式中作线性换元 \(x=\sqrt2\,\varepsilon u\)，得到
\[
 \int_{\mathbb R}e^{-x^2/(2\varepsilon^2)}\,\mathrm dx
 =\sqrt2\,\varepsilon\int_{\mathbb R}e^{-u^2}\,\mathrm du
 =\sqrt{2\pi}\,\varepsilon.
\]
这说明 \(g_\varepsilon\) 可积且积分为一；其正性、连续性由公式直接给出。

多项式乘 Gaussian 的绝对可积性来自指数衰减；
例如在充分大的 \(|x|\) 上，
\(|x|^k e^{-x^2/2}\leq C_k e^{-|x|/4}\)，其中 \(k=1,2\)。
奇偶性给出 \(\int x\cdot g_1(x)\,\mathrm dx=0\)。
又有 \(g_1'(x)=-x\cdot g_1(x)\)，故在 \([-R,R]\) 上分部积分，令 \(R\to\infty\)，得到
\[
 \int_{\mathbb R}x^2\cdot g_1(x)\,\mathrm dx
 =-\bigl[x\cdot g_1(x)\bigr]_{-\infty}^{+\infty}
     +\int_{\mathbb R}g_1(x)\,\mathrm dx
 =1.
\]
这里边界项趋于零，各积分绝对收敛。
换元 \(x=\varepsilon u\) 即给出一般 \(\varepsilon\) 的两个矩公式。

令
\[
 I(t)=\int_{\mathbb R}e^{itx}\cdot g_1(x)\,\mathrm dx.
\]
实参数差商满足
\[
 \left|\frac{e^{i(t+h)x}-e^{itx}}{h}\right|\leq|x|
 \quad(h\ne0).
\]
由于 \(|x|g_1(x)\) 可积，控制收敛允许对实部、虚部在积分号下求导：
\[
 I'(t)=i\int_{\mathbb R}x\cdot e^{itx}\cdot g_1(x)\,\mathrm dx.
\]
再在 \([-R,R]\) 上分部积分，并利用 \(g_1'=-x\cdot g_1\)，得到
\[
 \int_{\mathbb R}x\cdot e^{itx}\cdot g_1(x)\,\mathrm dx
 =-\bigl[e^{itx}\cdot g_1(x)\bigr]_{-\infty}^{+\infty}
       +it\int_{\mathbb R}e^{itx}\cdot g_1(x)\,\mathrm dx
 =it\cdot I(t).
\]
因而 \(I'(t)=-tI(t)\)，且 \(I(0)=1\)。
函数 \(e^{t^2/2}\cdot I(t)\) 的导数为零，于是
\(I(t)=e^{-t^2/2}\)。
用 \(x=\varepsilon u\) 换元便得到\eqref{b1:eq:30-gaussian-characteristic}。

最后，以 \(s=\varepsilon t\) 换元，再次使用刚求出的 \(I\)，可得
\[
 \int_{\mathbb R}e^{-itx}\cdot e^{-\varepsilon^2t^2/2}\,\mathrm dt
 =\frac{\sqrt{2\pi}}{\varepsilon}\cdot I(-x/\varepsilon)
 =\frac{\sqrt{2\pi}}{\varepsilon}
       \exp\left(-\frac{x^2}{2\varepsilon^2}\right).
\]
除以 \(2\pi\)，即是\eqref{b1:eq:30-gaussian-inversion}。
其中指数因子的模为一，剩余 Gaussian 因子可积，
所以最后的积分同样绝对收敛。
\end{proof}

\begin{definition}\label{b1:def:30-normal-distribution}
设 \(m\in\mathbb R\)、\(\sigma\geq0\)。
当 \(\sigma>0\) 时，把具有密度
\[
 x\longmapsto g_\sigma(x-m)
 =\frac{1}{\sqrt{2\pi}\,\sigma}
       \exp\left(-\frac{(x-m)^2}{2\sigma^2}\right)
\]
的 Borel 概率分布称为均值 \(m\)、方差 \(\sigma^2\) 的
\term[zhengtai-fenbu]{正态分布}[normal distribution]，
记为 \(N(m,\sigma^2)\)。
约定 \(N(m,0)=\delta_m\)。
其中 \(N(0,1)\) 称为
\term[biaozhun-zhengtai-fenbu]{标准正态分布}[standard normal distribution]，
其分布函数记为
\[
 \Phi(x)=\frac{1}{\sqrt{2\pi}}\int_{-\infty}^x e^{-u^2/2}\,\mathrm du.
\]
\end{definition}
\symidx{\(N(m,\sigma^2)\)：均值 \(m\)、方差 \(\sigma^2\) 的正态分布}
\symidx{\(\Phi\)：标准正态分布函数}

若 \(G\) 的分布为 \(N(0,1)\)，则 \(m+\sigma G\) 的分布为
\(N(m,\sigma^2)\)；当 \(\sigma>0\) 时由密度换元得到，
当 \(\sigma=0\) 时直接得到单点分布。
Gaussian 引理同时验证了名称中的均值与方差，并给出
\begin{equation}\label{b1:eq:30-normal-characteristic}
 \varphi_{N(m,\sigma^2)}(t)
 =\exp\left(imt-\frac{\sigma^2t^2}{2}\right).
\end{equation}
\(\Phi\) 连续且严格递增，满足
\(\Phi(-x)=1-\Phi(x)\)，在两端的极限分别为 \(0\)、\(1\)。
这些性质都来自正密度 \(g_1\) 的连续性、对称性及积分为一。

\subsection{特征函数的唯一性}
\label{b1:subsec:300602}

两个分布若对所有复指数函数给出相同积分，为什么必定相等？
我们先给分布加上一层 Gaussian 平滑。
平滑后的密度可以由特征函数的一个绝对收敛积分写出；
随后让平滑尺度趋于零，就能恢复原来的分布。

\begin{lemma}\label{b1:lem:30-gaussian-smoothing}
设 \(\mu\) 是实轴上的 Borel 概率测度，\(\varepsilon>0\)，并令
\[
 p_{\mu,\varepsilon}(x)
   =\int_{\mathbb R}g_\varepsilon(x-y)\,\mathrm d\mu(y),
 \quad x\in\mathbb R.
\]
则 \(p_{\mu,\varepsilon}\) 是连续概率密度，且
\begin{equation}\label{b1:eq:30-smoothed-density}
 p_{\mu,\varepsilon}(x)
 =\frac{1}{2\pi}\int_{\mathbb R}
       e^{-itx}\cdot\varphi_\mu(t)\cdot e^{-\varepsilon^2t^2/2}\,\mathrm dt.
\end{equation}
右端绝对收敛。若
\(\rho_{\mu,\varepsilon}=p_{\mu,\varepsilon}(x)\,\mathrm dx\)，则
\(\rho_{\mu,\varepsilon}\Rightarrow\mu\) 当 \(\varepsilon\downarrow0\)。
\end{lemma}
\begin{proof}
被积函数非负，Tonelli 定理与 \(\int g_\varepsilon=1\) 给出
\[
 \int_{\mathbb R}p_{\mu,\varepsilon}(x)\,\mathrm dx
 =\int_{\mathbb R}\left(\int_{\mathbb R}
       g_\varepsilon(x-y)\,\mathrm dx\right)\mathrm d\mu(y)
 =1.
\]
当 \(x_n\to x\) 时，\(g_\varepsilon(x_n-y)\to g_\varepsilon(x-y)\)，
而这些函数都由常数 \((\sqrt{2\pi}\,\varepsilon)^{-1}\) 控制。
该常数对概率测度 \(\mu\) 可积，故控制收敛给出密度的连续性。

将\eqref{b1:eq:30-gaussian-inversion}代入定义，可写成
\[
 p_{\mu,\varepsilon}(x)
 =\frac{1}{2\pi}\int_{\mathbb R}\int_{\mathbb R}
       e^{-it(x-y)}\cdot e^{-\varepsilon^2t^2/2}
       \,\mathrm dt\,\mathrm d\mu(y).
\]
交换积分所需的绝对积分为
\[
 \int_{\mathbb R}\int_{\mathbb R}
       e^{-\varepsilon^2t^2/2}\,\mathrm dt\,\mathrm d\mu(y)
 =\int_{\mathbb R}e^{-\varepsilon^2t^2/2}\,\mathrm dt<\infty.
\]
因此 Fubini 定理适用，先对 \(y\) 积分便得到
\eqref{b1:eq:30-smoothed-density}。

为证明弱收敛，任取实值 \(f\in C_b(\mathbb R)\)。
Fubini 定理及换元 \(x=y+\varepsilon z\) 给出
\[
 \int_{\mathbb R}f\,\mathrm d\rho_{\mu,\varepsilon}
 =\int_{\mathbb R}\int_{\mathbb R}
       f(y+\varepsilon z)\cdot g_1(z)\,\mathrm dz\,\mathrm d\mu(y).
\]
可以把右端看作固定乘积概率测度
\(\mu(\mathrm dy)\otimes(g_1(z)\,\mathrm dz)\) 上的积分。
对每个 \(y,z\)，连续性给出 \(f(y+\varepsilon z)\to f(y)\)，
而其绝对值不超过 \(\|f\|_\infty\)。
控制收敛遂给出极限 \(\int f\,\mathrm d\mu\)。
这恰好是第12章的弱收敛定义；这里不要求 \(f\) 在整个实轴上一致连续。
\end{proof}

\begin{theorem}[特征函数的唯一性]\label{b1:thm:30-characteristic-uniqueness}
设 \(\mu,\nu\) 是 \(\mathbb R\) 上的 Borel 概率测度。
若对每个 \(t\in\mathbb R\) 都有 \(\varphi_\mu(t)=\varphi_\nu(t)\)，
则 \(\mu=\nu\)。
\end{theorem}
\begin{proof}
由\eqref{b1:eq:30-smoothed-density}，对每个 \(\varepsilon>0\)，
都有 \(p_{\mu,\varepsilon}=p_{\nu,\varepsilon}\)，从而
\(\rho_{\mu,\varepsilon}=\rho_{\nu,\varepsilon}\)。
对任意实值 \(f\in C_b(\mathbb R)\)，使用平滑引理令 \(\varepsilon\downarrow0\)，得到
\[
 \int_{\mathbb R}f\,\mathrm d\mu
 =\lim_{\varepsilon\downarrow0}\int_{\mathbb R}f\,\mathrm d\rho_{\mu,\varepsilon}
 =\lim_{\varepsilon\downarrow0}\int_{\mathbb R}f\,\mathrm d\rho_{\nu,\varepsilon}
 =\int_{\mathbb R}f\,\mathrm d\nu.
\]
实轴是度量空间，\(\mu,\nu\) 都是有限正 Borel 测度，
所以\cref{b1:prop:weak-mass-uniqueness}给出 \(\mu=\nu\)。
\end{proof}

因此，特征函数并未丢失分布信息。
例如，要证明一个独立和仍服从某个正态分布，
只需计算其特征函数，并与\eqref{b1:eq:30-normal-characteristic}比较。
但在处理一列分布时，除了辨认候选极限，还必须保证质量没有逃向无穷远。

\subsection{紧性与子列极限的辨认}
\label{b1:subsec:300603}

第12章的紧性要求同一个紧集控制全族的尾质量。
特征函数如何提供这样的控制？在原点附近，
\(1-\operatorname{Re}\varphi_\mu(t)\) 是 \(1-\cos(tx)\) 的平均值。
再对一个短参数区间求平均，便能避免单个余弦函数不断回到一的困难。
这一尾概率估计的同类形式见
Fremlin\cite[285J]{Fremlin2016Measure}；下面给出本节所用的常数和证明。

\begin{lemma}\label{b1:lem:30-characteristic-tail}
设 \(\mu\) 是实轴上的 Borel 概率测度，\(\delta>0\)。
则
\begin{equation}\label{b1:eq:30-characteristic-tail}
 \mu\bigl(\{x:|x|\geq2/\delta\}\bigr)
 \leq\frac{1}{\delta}\int_{-\delta}^{\delta}
       \bigl(1-\operatorname{Re}\varphi_\mu(t)\bigr)\,\mathrm dt.
\end{equation}
\end{lemma}
\begin{proof}
函数 \(1-\cos(tx)\) 非负且不超过 \(2\)，
在 \([-\delta,\delta]\times\mathbb R\) 上对相应乘积测度可积。
交换积分并计算内层积分，得到
\[
 \frac{1}{2\delta}\int_{-\delta}^{\delta}
       \bigl(1-\operatorname{Re}\varphi_\mu(t)\bigr)\,\mathrm dt
 =\int_{\mathbb R}\left(1-\frac{\sin(\delta x)}{\delta x}\right)\mathrm d\mu(x),
\]
其中在 \(x=0\) 处将比值取为连续值 \(1\)。
右端被积函数处处非负，因为它也是 \(1-\cos(tx)\) 的区间平均。
当 \(|x|\geq2/\delta\) 时，
\(\sin(\delta x)/(\delta x)\leq1/2\)，故被积函数至少为 \(1/2\)。
于是右端至少为
\(\frac12\cdot\mu(\{x:|x|\geq2/\delta\})\)，这就给出所需不等式。
\end{proof}

\begin{lemma}\label{b1:lem:30-characteristic-tightness}
设 \((\mu_n)_{n\geq1}\) 是实轴上的 Borel 概率测度列。
若对每个 \(t\in\mathbb R\)，
\(\varphi_{\mu_n}(t)\to\psi(t)\)，且函数 \(\psi:\mathbb R\to\mathbb C\)
在 \(0\) 处连续，则 \(\{\mu_n:n\geq1\}\) 是紧族。
\end{lemma}
\begin{proof}
因为各 \(\varphi_{\mu_n}(0)=1\)，所以 \(\psi(0)=1\)。
给定 \(\eta>0\)，由零点处连续性，取 \(\delta>0\)，使
\[
 |1-\psi(t)|<\eta/8\quad(|t|\leq\delta).
\]
极限函数可测，且 \(|\psi(t)|\leq1\)。
在有限区间 \([-\delta,\delta]\) 上，
\(0\leq1-\operatorname{Re}\varphi_{\mu_n}(t)\leq2\)，
故控制收敛给出
\[
 \frac{1}{\delta}\int_{-\delta}^{\delta}
       \bigl(1-\operatorname{Re}\varphi_{\mu_n}(t)\bigr)\,\mathrm dt
 \longrightarrow
 \frac{1}{\delta}\int_{-\delta}^{\delta}
       \bigl(1-\operatorname{Re}\psi(t)\bigr)\,\mathrm dt
 \leq\eta/4.
\]
因此存在 \(N\)，使所有 \(n\geq N\) 的左端都小于 \(\eta/2\)。
由尾概率估计，这些 \(\mu_n\) 在
\(\{|x|\geq2/\delta\}\) 上的质量都小于 \(\eta/2\)。

还须把前面有限个测度纳入同一个控制。
对每个 \(n<N\)，由
\([-R,R]\uparrow\mathbb R\) 及测度连续性，可取足够大的 \(R_n\)，
使 \(\mu_n(\mathbb R\setminus[-R_n,R_n])<\eta\)。
取 \(R\) 不小于 \(2/\delta\) 及这些有限个 \(R_n\)，便有
\[
 \sup_{n\geq1}\mu_n(\mathbb R\setminus[-R,R])<\eta.
\]
闭区间 \([-R,R]\) 紧，所以满足\cref{b1:def:tight-family}的要求。
\end{proof}

\begin{proposition}\label{b1:prop:30-tight-characteristic-limit}
设 \((\mu_n)_{n\geq1}\) 是实轴上的 Borel 概率测度列，
\(\mu\) 是实轴上的 Borel 概率测度。
若 \(\{\mu_n:n\geq1\}\) 为紧族，且
\(\varphi_{\mu_n}(t)\to\varphi_\mu(t)\) 对每个 \(t\in\mathbb R\) 成立，
则 \(\mu_n\Rightarrow\mu\)。
\end{proposition}
\begin{proof}
实轴在通常度量下完备且可分，因而是波兰空间。
由\cref{b1:thm:prokhorov}，任何子列都有一个弱收敛的子子列；
设其中一个子子列的极限为 Borel 概率测度 \(\nu\)。
对每个固定 \(t\)，函数 \(\cos(tx)\)、\(\sin(tx)\) 都属于
\(C_b(\mathbb R)\)，所以弱收敛给出该子子列的特征函数趋于
\(\varphi_\nu(t)\)。
原来的点态极限假设又给出 \(\varphi_\nu(t)=\varphi_\mu(t)\)。
唯一性定理于是说明 \(\nu=\mu\)。
所有弱收敛子列的极限都被辨认为同一个 \(\mu\)，
再由\cref{b1:cor:tight-unique-limit}得到整列弱收敛。
\end{proof}

\subsection{特征函数的连续性定理}
\label{b1:subsec:300604}

前一命题已经假定候选概率测度 \(\mu\) 存在。
Fremlin\cite[285L--M]{Fremlin2016Measure}讨论了给定极限测度后的弱收敛判别和唯一性。
对中心极限定理而言，我们最先算出的却只是一个函数极限。
下面的\term[levy-lianxuxing-dingli]{Lévy 连续性定理}[Lévy's continuity theorem]
同时解决极限分布的存在与辨认。

\begin{theorem}[Lévy 连续性定理]\label{b1:thm:30-levy-continuity}
设 \((\mu_n)_{n\geq1}\) 是 \(\mathbb R\) 上的 Borel 概率测度列。
若存在函数 \(\psi:\mathbb R\to\mathbb C\)，使
\[
 \varphi_{\mu_n}(t)\longrightarrow\psi(t)
 \quad\text{对每个 }t\in\mathbb R,
\]
且 \(\psi\) 在 \(0\) 处连续，则存在唯一的 Borel 概率测度 \(\mu\)，
满足 \(\varphi_\mu=\psi\)，并且 \(\mu_n\Rightarrow\mu\)。
反过来，若 \(\mu_n\Rightarrow\mu\)，其中 \(\mu\) 是实轴上的
Borel 概率测度，则
\(\varphi_{\mu_n}(t)\to\varphi_\mu(t)\) 对每个实数 \(t\) 成立，
且极限特征函数在 \(0\) 处连续。
\end{theorem}
\begin{proof}
先证正向结论。\cref{b1:lem:30-characteristic-tightness}说明测度列为紧族。
由于 \(\mathbb R\) 是波兰空间，Prokhorov 定理给出一个弱收敛子列
\(\mu_{n_k}\Rightarrow\mu\)，其极限仍是 Borel 概率测度。
用有界连续测试函数 \(\cos(tx)\)、\(\sin(tx)\) 分别取极限，得到
\[
 \varphi_\mu(t)=\lim_{k\to\infty}\varphi_{\mu_{n_k}}(t)=\psi(t)
 \quad(t\in\mathbb R).
\]
这证明了具有所求特征函数的概率测度确实存在。
唯一性由\cref{b1:thm:30-characteristic-uniqueness}给出；
再用\cref{b1:prop:30-tight-characteristic-limit}，便得到整列
\(\mu_n\Rightarrow\mu\)。

反向结论直接用同样的余弦、正弦测试函数。
所得极限是 \(\varphi_\mu\)，它由
\cref{b1:prop:30-characteristic-basic}在整个实轴上一致连续，
特别在 \(0\) 处连续。
\end{proof}

特别地，当候选概率测度 \(\mu\) 已经给定时，
\(\mu_n\Rightarrow\mu\) 等价于所有参数 \(t\) 下的特征函数收敛。
此时无需另行检查紧性，因为极限 \(\varphi_\mu\) 在零点连续，
已经通过尾概率估计保证紧性。

\ProseParagraphBreak

零点连续性不能删去。
取 \(\mu_n=N(0,n^2)\)，则
\[
 \varphi_{\mu_n}(t)=e^{-n^2t^2/2}
 \longrightarrow
 \begin{cases}
  1,&t=0,\\
  0,&t\ne0.
 \end{cases}
\]
极限在零点不连续，因而不是任何概率测度的特征函数。
从密度也能直接看见质量逃逸：对每个固定 \(R>0\)，
\[
 \mu_n([-R,R])
 \leq\frac{2R}{\sqrt{2\pi}\,n}\longrightarrow0.
\]
任何固定紧集都装不下这列分布的大部分质量。
连续性定理中的零点条件正是把这种情形排除的条件。
下一节将计算标准化独立和的特征函数极限；
一旦极限被认作标准正态特征函数，本节的定理便把函数极限变成分布极限。

% !TeX root = ../../Book1.tex
% 实读来源：Fremlin, Measure Theory, Vol. 2，作者在线第27章摘本
% 274A、274H(c)、274I，及第28章285F--G。正式引文按稳定条号定位。
% 第27章摘本 PDF 第36--37页（274A）、第43页（274H(c)）、第44页（274I）。
% Fremlin274I经Lindeberg条件调用274G；274H(c)说明其主线不使用特征函数。
% 以下采用本章自证的特征函数路线。
% Bass, Real Analysis for Graduate Students, v5，印刷第305--307页，
% 仅作教学背景。复数幂极限与下述无限方差模型均在正文给出完整推导。
\section{独立同分布变量的中心极限定理}
\label{b1:sec:3007}

设独立同分布变量的均值为 \(m\)，前两节的大数律说明
\((X_1+\cdots+X_n)/n\) 趋于 \(m\)。
这个结论没有描述平均误差的形状。
当方差为有限正数 \(\sigma^2\) 时，独立和的方差是 \(n\sigma^2\)，
因此 \(\sigma\sqrt n\) 是使中心化和保持单位方差的尺度。
本节证明：在这个尺度下，极限分布只保留均值和方差信息，
而与原分布是离散还是连续无关。

\subsection{有限方差与特征函数的局部展开}
\label{b1:subsec:300701}

独立和的特征函数是各项特征函数的乘积。
标准化又把参数 \(t\) 缩成 \(t/\sqrt n\)，
所以只需知道原特征函数在零点附近的二阶行为。
关键是用二阶矩控制余项，不能在 Taylor 展开中暗中加入三阶矩条件。

\begin{proposition}\label{b1:prop:30-characteristic-second-order}
设 \(X\) 是概率空间 \((\Omega,\mathcal F,\mathbb P)\) 上的实随机变量，
且 \(\mathbb E(X^2)<\infty\)。
则特征函数 \(\varphi_X\) 关于实参数属于 \(C^2(\mathbb R;\mathbb C)\)，并满足
\[
 \varphi_X'(t)=i\mathbb E(X\cdot e^{itX}),
 \quad
 \varphi_X''(t)=-\mathbb E(X^2\cdot e^{itX}).
\]
当 \(t\to0\) 时，有
\begin{equation}\label{b1:eq:30-characteristic-expansion}
 \varphi_X(t)
 =1+it\mathbb EX-\frac{t^2}{2}\mathbb E(X^2)+o(t^2).
\end{equation}
特别地，若 \(\mathbb EX=0\)、\(\operatorname{Var}(X)=1\)，则
\(\varphi_X(t)=1-t^2/2+o(t^2)\)。
\end{proposition}
\begin{proof}
由 Cauchy--Schwarz 不等式及 \(\mathbb P(\Omega)=1\)，
\(\mathbb E|X|\leq(\mathbb E(X^2))^{1/2}<\infty\)。
对 \(h\ne0\)，有
\[
 \left|\frac{e^{i(t+h)X}-e^{itX}}{h}\right|\leq|X|.
\]
差商逐点趋于 \(iX\cdot e^{itX}\)，故控制收敛给出一阶导数公式。
再对该公式求差商，其被积函数的模不超过 \(X^2\)，
控制收敛给出二阶导数公式。
当 \(t_j\to t\) 时，两个导数公式中的被积函数分别由
\(|X|\)、\(X^2\) 控制，故两个导数也都连续。

为了明确二阶余项所需的控制，对每个实数 \(u\)，两次使用微积分基本定理可得
\[
 e^{iu}
 =1+iu-u^2\int_0^1(1-s)e^{isu}\,\mathrm ds
 =1+iu-\frac{u^2}{2}+u^2\cdot r(u),
\]
其中
\[
 r(u)=-\int_0^1(1-s)(e^{isu}-1)\,\mathrm ds.
\]
这里 \(r(0)=0\)，\(|r(u)|\leq1\)，且 \(r(u)\to0\) 当 \(u\to0\)。
这些性质分别来自定义、\(|e^{isu}-1|\leq2\) 及区间上的控制收敛。
这个积分余项对正、负的 \(u\) 都成立，因为积分变量始终是 \(s\in[0,1]\)。

将 \(u=tX\) 代入并取期望，得到
\[
 \varphi_X(t)
 =1+it\mathbb EX-\frac{t^2}{2}\mathbb E(X^2)
       +t^2\mathbb E\bigl(X^2\cdot r(tX)\bigr).
\]
由于 \(X\) 为有限实值随机变量，\(r(tX)\to0\) 逐点成立；
而
\[
 |X^2\cdot r(tX)|\leq X^2
\]
给出与 \(t\) 无关的可积控制。
因此最后一个期望趋于零，证明\eqref{b1:eq:30-characteristic-expansion}。
在均值零、方差一的情形，\(\mathbb E(X^2)=1\)，即得最后的展开。
\end{proof}

若改用余项的三次幂上界，再对 \(X\) 积分，就会无端要求
\(\mathbb E|X|^3<\infty\)。
上面的证明保留了有界函数 \(r\)，从而只需二阶矩。
这也是后面中心极限定理恰好适用于有限方差变量的原因。

\subsection{标准化和与 Lindeberg–Lévy 定理}
\label{b1:subsec:300702}

从单个特征函数的局部展开到独立和，还需要一个复数幂的极限。
虽然最终极限是正实数，原特征函数的值一般是复数；
下面直接比较整数次幂和指数，避免选择复对数的分支。

\begin{lemma}\label{b1:lem:30-complex-power-limit}
设 \((z_n)_{n\geq1}\) 是复数列，\(z\in\mathbb C\)。
若 \(n(z_n-1)\to z\)，则 \(z_n^n\to e^z\)。
\end{lemma}
\begin{proof}
记 \(w_n=z_n-1\)，则 \(n w_n\to z\)，所以 \(n|w_n|\) 有界。
这里复指数可以直接由绝对收敛级数
\(e^w=\sum_{k=0}^\infty w^k/k!\) 定义。
其乘法公式由绝对收敛级数的乘积和二项式公式得到，
因而 \((e^w)^n=e^{nw}\)。
级数还给出
\[
 |e^w-1-w|
 \leq\sum_{k=2}^\infty\frac{|w|^k}{k!}
 \leq\frac{|w|^2}{2}\cdot e^{|w|},
 \quad
 |1+w|\leq e^{|w|},\quad |e^w|\leq e^{|w|}.
\]
对 \(a=1+w_n\)、\(b=e^{w_n}\) 使用
\(a^n-b^n=(a-b)\sum_{j=0}^{n-1}a^{n-1-j}b^j\)，得到
\[
 |(1+w_n)^n-e^{n w_n}|
 \leq\frac n2\cdot|w_n|^2\cdot e^{n|w_n|}.
\]
其中 \(e^{n|w_n|}\) 有界，且
\[
 n|w_n|^2=\frac{(n|w_n|)^2}{n}\longrightarrow0.
\]
另一方面，指数函数的连续性给出 \(e^{n w_n}\to e^z\)，故结论成立。
\end{proof}

下面的结果称为独立同分布变量的
\term[zhongxin-jixian-dingli]{中心极限定理}[central limit theorem]，
这一有限方差版本也称
\term[lindeberg-levy-dingli]{Lindeberg--Lévy 定理}[Lindeberg--Lévy theorem]。
其陈述可对照 Fremlin\cite[274I]{Fremlin2016Measure}；
这里采用前一节建立的特征函数方法证明。

\begin{theorem}[Lindeberg--Lévy]\label{b1:thm:30-iid-clt}
设 \((X_k)_{k\geq1}\) 是概率空间
\((\Omega,\mathcal F,\mathbb P)\) 上的独立同分布实随机变量列，且
\[
 \mathbb E(X_1^2)<\infty,\quad
 m=\mathbb EX_1,\quad
 \sigma^2=\operatorname{Var}(X_1)>0.
\]
令 \(S_n=\sum_{k=1}^nX_k\)，并定义
\[
 Z_n=\frac{S_n-nm}{\sigma\sqrt n}.
\]
则 \(Z_n\) 依分布收敛于 \(N(0,1)\)；等价地，对每个 \(x\in\mathbb R\)，
\begin{equation}\label{b1:eq:30-clt-cdf}
 \mathbb P\left(\frac{S_n-nm}{\sigma\sqrt n}\leq x\right)
 \longrightarrow\Phi(x).
\end{equation}
\end{theorem}
\begin{proof}
令 \(Y_k=(X_k-m)/\sigma\)。
可测函数保持独立性，且这些 \(Y_k\) 同分布，
均满足 \(\mathbb EY_k=0\)、\(\mathbb E(Y_k^2)=1\)。
记公共特征函数为 \(\chi\)。
由\cref{b1:prop:30-characteristic-second-order}，
存在在零点趋于零的函数 \(\eta\)，使
\[
 \chi(u)=1-\frac{u^2}{2}+u^2\cdot\eta(u).
\]
独立和及伸缩的特征函数公式给出
\[
 \varphi_{Z_n}(t)=\left(\chi(t/\sqrt n)\right)^n.
\]
固定 \(t\ne0\)，则
\[
 n\bigl(\chi(t/\sqrt n)-1\bigr)
 =-\frac{t^2}{2}+t^2\cdot\eta(t/\sqrt n)
 \longrightarrow-\frac{t^2}{2}.
\]
由复数幂极限引理，
\(\varphi_{Z_n}(t)\to e^{-t^2/2}\)。
对 \(t=0\)，两端恒等于一，所以点态收敛在所有实数参数上成立。

由\eqref{b1:eq:30-normal-characteristic}，极限函数是
\(N(0,1)\) 的特征函数，且在零点连续。
\cref{b1:thm:30-levy-continuity}于是给出 \(Z_n\) 的分布弱收敛于
\(N(0,1)\)。
最后，\(\Phi\) 在每个实数点连续，
故\cref{b1:prop:cdf-weak-convergence}将弱收敛化为
\eqref{b1:eq:30-clt-cdf}。
\end{proof}

中心化消去特征函数的一阶项，平方根尺度使每个因子的二阶项为
\(-t^2/(2n)\)，而 \(n\) 个独立因子的乘积把这一项累积成
\(e^{-t^2/2}\)。
高阶细节在极限中消失；在本证明中，这件事由有限二阶矩下的受控余项完成。

\subsection{退化情形与具体分布模型}
\label{b1:subsec:300703}

若方差为零，就不能除以 \(\sigma\)。
这不是另一种波动分布，而是根本没有随机波动。
对独立同分布列，若
\(\operatorname{Var}(X_1)=\mathbb E((X_1-m)^2)=0\)，
则对每个 \(k\)，都有 \(X_k=m\) 几乎必然。
这些概率一事件的可数交仍有概率一，
所以在同一个概率一集合上，所有 \(n\) 都满足 \(S_n=nm\)。

\begin{corollary}\label{b1:cor:30-clt-degenerate}
设 \((X_k)_{k\geq1}\) 是概率空间
\((\Omega,\mathcal F,\mathbb P)\) 上的独立同分布实随机变量列，
\(\mathbb E(X_1^2)<\infty\)，并令
\[
 m=\mathbb EX_1,\quad \sigma^2=\operatorname{Var}(X_1),
 \quad S_n=\sum_{k=1}^nX_k.
\]
允许 \(\sigma=0\)，则
\[
 \frac{S_n-nm}{\sqrt n}\ \Rightarrow\ N(0,\sigma^2).
\]
\end{corollary}
\begin{proof}
若 \(\sigma=0\)，由平方积分为零及可数交的论证，
左端对每个 \(n\) 几乎必然为零，其分布为
\(\delta_0=N(0,0)\)。
若 \(\sigma>0\)，记中心极限定理中的标准化和为 \(Z_n\)。
对任意实值 \(f\in C_b(\mathbb R)\)，函数 \(u\mapsto f(\sigma u)\)
仍有界连续，因此 \(Z_n\Rightarrow N(0,1)\) 给出
\[
 \mathbb E f(\sigma Z_n)
 \longrightarrow\int_{\mathbb R}f(\sigma u)\cdot g_1(u)\,\mathrm du.
\]
右端正是 \(f\) 对 \(N(0,\sigma^2)\) 的积分，故得到所需弱收敛。
\end{proof}

正态分布给出一个每一步都能精确计算的模型。
设 \(X_k\) 独立同分布于 \(N(m,\sigma^2)\)，其中 \(\sigma>0\)。
对每个正整数 \(n\)，独立性给出
\[
 \varphi_{S_n}(t)
 =\exp\left(inmt-\frac{n\sigma^2t^2}{2}\right).
\]
特征函数唯一性说明 \(S_n\) 的分布恰好是 \(N(nm,n\sigma^2)\)，
所以 \((S_n-nm)/(\sigma\sqrt n)\) 对每个 \(n\) 已经服从 \(N(0,1)\)。
中心极限定理则说明，原来的变量无需具有正态分布，
在极限中仍会出现这个模型。

\ProseParagraphBreak

先看离散模型。设 \(0<p<1\)，\(X_k\) 独立且具有同一 Bernoulli 分布，
即 \(\mathbb P(X_k=1)=p\)、\(\mathbb P(X_k=0)=1-p\)。
令 \(K_n=\sum_{k=1}^nX_k\)，它记录前 \(n\) 次试验中取值一的次数。
对每个 \(r=0,\ldots,n\)，指定一个恰有 \(r\) 项取值一的模式，
其概率由独立性等于 \(p^r\cdot(1-p)^{n-r}\)。
这样的互不相交模式共有 \(\binom nr\) 个，故
\[
 \mathbb P(K_n=r)=\binom nr p^r\cdot(1-p)^{n-r}.
\]
这个分布称为参数 \(n,p\) 的
\term[erxiang-fenbu]{二项分布}[binomial distribution]。
由 \(X_k^2=X_k\)，得到
\(\mathbb EX_k=p\)、\(\operatorname{Var}(X_k)=p(1-p)\)。
因此
\[
 \frac{K_n-np}{\sqrt{np(1-p)}}\ \Rightarrow\ N(0,1).
\]
对任意固定实数 \(a<b\)，将两个分布函数的极限相减，得到
\[
 \mathbb P\left(
 a<\frac{K_n-np}{\sqrt{np(1-p)}}\leq b
 \right)
 \longrightarrow\Phi(b)-\Phi(a).
\]
离散取值并不妨碍分布趋于连续分布。
若 \(p=0\) 或 \(p=1\)，方差为零，应使用前面的退化结论。

\ProseParagraphBreak

再看连续模型。设 \(a<b\)，\(X_k\) 独立且在 \([a,b]\) 上均匀分布。
直接按密度积分可得
\[
 m=\frac{a+b}{2},\quad
 \mathbb E((X_k-m)^2)
 =\frac{1}{b-a}\int_a^b(x-m)^2\,\mathrm dx
 =\frac{(b-a)^2}{12}.
\]
它的特征函数也可以显式写出：
\[
 \varphi_{X_k}(t)
 =e^{itm}\cdot\frac{\sin\bigl(t(b-a)/2\bigr)}{t(b-a)/2},
\]
其中 \(t=0\) 处取连续值一。
对 \(t\ne0\)，这个公式由
\((b-a)^{-1}\int_a^b e^{itx}\,\mathrm dx\) 直接积分得到。
中心极限定理给出
\[
 \frac{\sqrt{12}}{(b-a)\sqrt n}
       \left(\sum_{k=1}^nX_k-\frac{n(a+b)}{2}\right)
 \ \Rightarrow\ N(0,1).
\]
Bernoulli 与均匀分布的形状很不同，
但中心化及单位方差归一化以后，二阶展开给出同一个极限。

\subsection{方差条件、极限方式与反例}
\label{b1:subsec:300704}

中心极限定理控制有界连续测试函数的积分，
也控制极限分布的连续性集合的概率。
它本身没有给出有限 \(n\) 时的误差阶。
更不能把这一结论理解为任意可测事件的概率都已经接近。
例如在参数 \(0<p<1\) 的 Bernoulli 模型中，令
\[
 A_n=\left\{\frac{r-np}{\sqrt{np(1-p)}}:r=0,\ldots,n\right\}.
\]
标准化和的分布在有限集 \(A_n\) 上赋值一，
而标准正态分布在 \(A_n\) 上赋值零。
因此，对这些随 \(n\) 改变的集合，两个概率始终相差一。

\ProseParagraphBreak

依分布收敛也不自动成为同一个概率空间上的依概率收敛。
设 \(X_1,X_2,\ldots\) 是独立标准正态变量，并令
\[
 S_n=\sum_{k=1}^nX_k,\quad Z_n=S_n/\sqrt n.
\]
每个 \(Z_n\) 的分布都是 \(N(0,1)\)，所以它们的分布当然收敛于
\(N(0,1)\)。
但定义
\[
 U_n=\frac{S_n}{\sqrt n},\quad
 V_n=\frac{S_{2n}-S_n}{\sqrt n}.
\]
不交分块保持独立，故 \(U_n,V_n\) 独立，且都服从 \(N(0,1)\)。
于是
\[
 Z_{2n}-Z_n
 =\left(\frac{1}{\sqrt2}-1\right)U_n+\frac{1}{\sqrt2}V_n
 \quad\text{服从 }N(0,2-\sqrt2).
\]
这里分布的认定仍只需特征函数：两个独立标准正态变量的线性组合
\(cU_n+dV_n\) 的特征函数为 \(e^{-(c^2+d^2)t^2/2}\)，
唯一性定理给出方差 \(c^2+d^2\) 的正态分布。
因此对任意固定 \(\varepsilon>0\)，
\[
 \mathbb P(|Z_{2n}-Z_n|>\varepsilon)
 =2\left[1-\Phi\left(\frac{\varepsilon}{\sqrt{2-\sqrt2}}\right)\right]>0,
\]
且此值与 \(n\) 无关。
若存在实随机变量 \(Z\) 使 \(Z_n\to Z\) 依概率成立，则三角不等式给出
\[
 \mathbb P(|Z_{2n}-Z_n|>\varepsilon)
 \leq
 \mathbb P(|Z_{2n}-Z|>\varepsilon/2)
 +\mathbb P(|Z_n-Z|>\varepsilon/2)
 \longrightarrow0,
\]
产生矛盾。
所以这列标准化和不依概率收敛于任何实随机变量，也不几乎必然收敛。

\ProseParagraphBreak

最后考察二阶矩条件。只有有限一阶矩时，大数律仍然成立，
但 \(\sqrt n\) 尺度的正态极限不再得到保证。
令
\[
 c=\left(\sum_{k=1}^\infty k^{-3}\right)^{-1},\quad
 \mathbb P(X=k)=\mathbb P(X=-k)=\frac{c}{2k^3}
 \quad(k=1,2,\ldots).
\]
这些概率之和为一，并且
\[
 \mathbb E|X|=c\sum_{k=1}^\infty k^{-2}<\infty,\quad
 \mathbb EX=0,\quad
 \mathbb E(X^2)=c\sum_{k=1}^\infty k^{-1}=\infty.
\]
其中均值为零使用了对称性和绝对可积性。
由\cref{b1:cor:30-iid-construction}，可以在同一个概率空间上
构造具有这个分布的独立同分布列 \(X_1,X_2,\ldots\)。
公共特征函数为
\[
 \chi(u)=c\sum_{k=1}^\infty\frac{\cos(ku)}{k^3}.
\]
该级数绝对收敛，且 \(\chi\) 为实值连续函数，\(\chi(0)=1\)。

\ProseParagraphBreak

由余弦在零点的二阶展开，可取 \(\delta_0>0\)，使
\(1-\cos v\geq v^2/4\) 在 \(|v|\leq\delta_0\) 时成立。
固定 \(t\ne0\)，对充分大的 \(n\)，令
\[
 M_n=\left\lfloor\frac{\delta_0\sqrt n}{|t|}\right\rfloor\geq1.
\]
由于 \(1-\cos v\geq0\)，可以只保留前 \(M_n\) 项，得到
\begin{align*}
 n\bigl(1-\chi(t/\sqrt n)\bigr)
 &=cn\sum_{k=1}^\infty
       \frac{1-\cos(kt/\sqrt n)}{k^3}\\
 &\geq\frac{ct^2}{4}\cdot\sum_{k=1}^{M_n}\frac1k
 \longrightarrow\infty.
\end{align*}
另一方面，\(\chi(t/\sqrt n)\to1\)，并且 \(\chi\leq1\)，
所以对充分大的 \(n\)，有 \(0<\chi(t/\sqrt n)\leq1\)。
利用实数不等式 \(1-v\leq e^{-v}\)，可得
\[
 0\leq\bigl(\chi(t/\sqrt n)\bigr)^n
 \leq\exp\left(-n\bigl(1-\chi(t/\sqrt n)\bigr)\right)
 \longrightarrow0.
\]
若 \(S_n=\sum_{k=1}^nX_k\)，则 \(S_n/\sqrt n\) 的特征函数
在每个非零 \(t\) 处趋于零，在 \(t=0\) 处恒为一。
这个点态极限不在零点连续，故由 Lévy 连续性定理的必要方向，
\(S_n/\sqrt n\) 不存在概率分布意义下的弱极限。

\ProseParagraphBreak

这个例子说明，删去二阶矩条件后，不能对所有可积独立同分布列
仍坚持平方根归一化并期待正态极限；其他归一化尺度属于进一步理论。
大数律、中心极限定理和随机变量本身的收敛分别回答不同的问题：
平均值是否稳定，波动分布在何种尺度下稳定，
以及同一个概率空间中的样本值是否趋于某个随机变量。
应用时须先确定所问的问题，再选择相应的假设和极限方式。


\begin{chaptersummary}
概率空间是总质量为一的测度空间，随机变量是可测映射，其分布是推前测度。
独立性是联合分布的乘积结构；由有限维分布构造可数乘积以后，
给定的一列实分布可以在同一个概率空间中实现为独立随机变量。

\ProseParagraphBreak

Markov 和 Chebyshev 估计把积分控制化为尾概率控制。
Borel--Cantelli 引理进一步把可求和的概率变为几乎必然只有有限次例外，
并在独立条件下给出反向结论。
弱大数律只需控制单个时刻的平均误差；强大数律还要控制一整段部分和，
因此需要极大不等式、可和级数及归一化工具。
截断同时处理大值出现的次数和截断变量的方差，使有限一阶矩足以得到强大数律。

\ProseParagraphBreak

特征函数把独立和的分布转成乘积。Gaussian 平滑恢复分布，
特征函数在原点附近的控制保证紧性，而紧性与唯一性共同识别弱极限。
有限二阶矩给出原点处的二阶展开；中心化后的一阶项消失，
平方根归一化使二阶项累积为 Gaussian 特征函数，得到中心极限定理。

\ProseParagraphBreak

大数律与中心极限定理采用不同的尺度：前者把和除以 \(n\)，
后者把中心化的和除以 \(\sqrt n\)。前者的可积性条件与后者的有限方差条件不能混用；
分布极限也不能自动升级为同一概率空间上的依概率或几乎必然极限。
\end{chaptersummary}

\BookExercises

% 自拟：把依概率收敛变为函数空间上的完备度量，训练代表与统一零测集的处理。
\begin{optionalexercise}\label{b1:ex:30-probability-metric}
设 \((\Omega,\mathcal F,\mathbb P)\) 为概率空间，\(L^0\) 为有限实值随机变量按几乎必然相等所得的等价类空间。
证明
\[
 \mathrm d(X,Y)=\mathbb E\min\{1,|X-Y|\}
\]
是 \(L^0\) 上的完备度量，并且 \(\mathrm d(X_n,X)\to0\) 当且仅当 \(X_n\) 依概率收敛于 \(X\)。
\end{optionalexercise}
\begin{solution}
几乎必然相等的代表给出相同积分，所以 \(\mathrm d\) 良定义。
它非负、对称，且为零当且仅当 \(X=Y\) 几乎必然。
由三角不等式及
\(\min(1,a+b)\leq\min(1,a)+\min(1,b)\)（\(a,b\geq0\)），得到度量的三角不等式。
对 \(0<\varepsilon<1\)，有
\[
 \varepsilon\cdot\mathbb P(|X-Y|>\varepsilon)
 \leq \mathrm d(X,Y)
 \leq\varepsilon+\mathbb P(|X-Y|>\varepsilon).
\]
左式说明度量收敛蕴含依概率收敛；右式先取上极限再令 \(\varepsilon\downarrow0\)，得到反向结论。

设 \((X_n)\) 关于 \(\mathrm d\) 为 Cauchy 列。
可取严格递增指标 \(n_k\)，使
\(\mathrm d(X_{n_{k+1}},X_{n_k})\leq2^{-2k}\)。上式给出
\[
 \sum_{k=1}^{\infty}\mathbb P(|X_{n_{k+1}}-X_{n_k}|>2^{-k})
 \leq\sum_{k=1}^{\infty}2^{-k}<\infty.
\]
由第一 Borel--Cantelli 引理，在一个概率为一的可测集上，
充分靠后的相邻差均不超过 \(2^{-k}\)，因此 \((X_{n_k})\) 逐点为实数 Cauchy 列。
在这个集合上取极限，在其补集上取零，得到有限实值可测函数 \(X\)。
控制收敛给出 \(\mathrm d(X_{n_k},X)\to0\)。最后，对充分大的 \(n\) 和 \(k\)，
\[
 \mathrm d(X_n,X)\leq \mathrm d(X_n,X_{n_k})+\mathrm d(X_{n_k},X),
\]
两项均可任意小，故全列收敛。这样得到完备性，而无须假设原概率空间已经完备化。
\end{solution}

% 自拟：以独立乘积期望和Markov估计导出Bernoulli平均值的指数尾界。
\begin{optionalexercise}\label{b1:ex:30-bernoulli-exponential-tail}
设 \((\Omega,\mathcal F,\mathbb P)\) 为概率空间，\((X_k)_{k\geq1}\) 为独立随机变量列，
且 \(\mathbb P(X_k=1)=p\)、\(\mathbb P(X_k=0)=1-p\)，其中 \(0<p<1\)。令 \(S_n=\sum_{k=1}^nX_k\)。
对 \(0<q<1\)，记
\[
 I_p(q)=q\log\frac qp+(1-q)\log\frac{1-q}{1-p}.
\]
证明：若 \(p<q<1\)，则 \(\mathbb P(S_n\geq nq)\leq e^{-nI_p(q)}\)；
若 \(0<q<p\)，则 \(\mathbb P(S_n\leq nq)\leq e^{-nI_p(q)}\)。
由此直接推出 \(S_n/n\to p\) 几乎必然。
\end{optionalexercise}
\begin{solution}
独立性给出 \(\mathbb E e^{tS_n}=(1-p+pe^t)^n\)。当 \(t>0\) 时，Markov 估计说明
\[
 \mathbb P(S_n\geq nq)
 \leq e^{-tnq}(1-p+pe^t)^n.
\]
若 \(q>p\)，取
\(t=\log\bigl(q(1-p)/(p(1-q))\bigr)>0\)，则
\(1-p+pe^t=(1-p)/(1-q)\)。代入并取对数，右端恰为 \(e^{-nI_p(q)}\)。
若 \(q<p\)，同一选择给出 \(t<0\)；此时事件 \(\{S_n\leq nq\}\)
蕴含 \(e^{tS_n}\geq e^{tnq}\)，用相同的 Markov 估计便得到下尾界。

有
\[
 I_p(p)=0,\quad I_p'(p)=0,\quad
 I_p''(q)=\frac1q+\frac1{1-q}>0.
\]
因此 \(q\ne p\) 时 \(I_p(q)>0\)。对任意
\(0<\varepsilon<\min(p,1-p)\)，两侧尾界给出
\[
 \sum_{n=1}^{\infty}\mathbb P(|S_n/n-p|\geq\varepsilon)<\infty.
\]
第一 Borel--Cantelli 引理说明固定误差最终消失。
再对一列趋于零的正有理误差取概率为一的可数交，得到几乎必然收敛。
\end{solution}

% 自拟：最大值的对数尺度，用稀疏子列弥补直接概率界不易求和的困难。
\begin{optionalexercise}\label{b1:ex:30-exponential-maximum}
设 \((\Omega,\mathcal F,\mathbb P)\) 为概率空间，\((X_k)_{k\geq1}\) 独立同分布，
各变量非负且 \(\mathbb P(X_k>x)=e^{-x}\)（\(x\geq0\)）。令 \(M_n=\max_{1\leq k\leq n}X_k\)。
证明
\[
 \frac{M_n}{\log n}\longrightarrow1\quad\text{几乎必然， }n\to\infty.
\]
\end{optionalexercise}
\begin{solution}
固定 \(0<\varepsilon<1\)。由独立性及 \(1-u\leq e^{-u}\)，
\[
 \mathbb P\bigl(M_n\leq(1-\varepsilon)\log n\bigr)
 =\bigl(1-n^{-(1-\varepsilon)}\bigr)^n
 \leq e^{-n^\varepsilon}.
\]
右侧关于 \(n\) 可和，因为充分大的 \(n\) 满足 \(n^\varepsilon\geq2\log n\)。
第一 Borel--Cantelli 引理给出几乎必然的下极限不小于 \(1-\varepsilon\)。

上尾用并集估计：
\[
 \mathbb P\bigl(M_n>(1+\varepsilon)\log n\bigr)
 \leq n\cdot n^{-(1+\varepsilon)}=n^{-\varepsilon}.
\]
这个界未必可和，但在 \(n=2^j\) 上可和。因此几乎必然地，对充分大的 \(j\)，
\(M_{2^j}\leq(1+\varepsilon)j\log2\)。
若 \(2^j\leq n<2^{j+1}\)，由最大值的单调性，
\[
 \frac{M_n}{\log n}
 \leq\frac{M_{2^{j+1}}}{j\log2}
 \leq(1+\varepsilon)\frac{j+1}{j}.
\]
于是上极限不超过 \(1+\varepsilon\)。最后在 \(\varepsilon\downarrow0\) 的可数列上同时使用上下界，即得结论。
\end{solution}

% 自拟：把可数多个强大数律和分布函数的单调性组合成一致估计。
\begin{optionalexercise}\label{b1:ex:30-empirical-cdf}
设 \((\Omega,\mathcal F,\mathbb P)\) 为概率空间，\((X_k)_{k\geq1}\) 为独立同分布的实随机变量，
共同分布函数 \(F\) 连续。定义\term[jingyan-fenbu-hanshu]{经验分布函数}[empirical distribution function]
\[
 F_n(x)=\frac1n\sum_{k=1}^n\mathbf1_{\{X_k\leq x\}}.
\]
证明 \(\sup_{x\in\mathbb R}|F_n(x)-F(x)|\to0\) 几乎必然。
\end{optionalexercise}
\begin{solution}
因为 \(F_n\) 右连续且 \(F\) 连续，所求上确界等于对有理数 \(x\) 所取的上确界，
因而是可测随机变量。
固定整数 \(m\geq2\)。由 \(F\) 连续及两端极限为 \(0,1\)，
可取 \(a_1<\cdots<a_{m-1}\)，使 \(F(a_j)=j/m\)。
对每个 \(a_j\)，变量列 \(\mathbf1_{\{X_k\leq a_j\}}\) 独立同分布且有界，
故强大数律给出 \(F_n(a_j)\to F(a_j)\) 几乎必然。
所有 \(m\) 和相应的有限个 \(j\) 合起来只有可数个条件，
因此存在同一个概率为一的集合，使这些收敛全部成立。

令 \(a_0=-\infty\)、\(a_m=+\infty\)，并把两个分布函数在这两个形式端点的值记为 \(0,1\)。
若 \(a_j\leq x\leq a_{j+1}\)，单调性给出
\[
 F_n(a_j)\leq F_n(x)\leq F_n(a_{j+1}),\quad
 F(a_j)\leq F(x)\leq F(a_{j+1}).
\]
由于相邻 \(F\) 值之差为 \(1/m\)，可得
\[
 \sup_{x\in\mathbb R}|F_n(x)-F(x)|
 \leq\max_{0\leq j\leq m}|F_n(a_j)-F(a_j)|+\frac1m.
\]
在上述概率为一的集合内先令 \(n\to\infty\)，再令 \(m\to\infty\)，便得到所需一致收敛。
有限网格上的误差连同网格间距，统一控制了整个实轴上的误差。
\end{solution}

% 自拟：用复随机变量的模长等号识别算术支集；不借助一般Fourier反演。
\begin{optionalexercise}\label{b1:ex:30-lattice-characteristic}
设 \((\Omega,\mathcal F,\mathbb P)\) 为概率空间，\(X\) 为实随机变量，
\(\phi(t)=\mathbb E e^{\mathrm itX}\)，\(t_0\ne0\)。证明
\[
 |\phi(t_0)|=1
 \quad\Longleftrightarrow\quad
 \text{存在 }a\in\mathbb R\text{，使 }
 \mathbb P\left(X\in a+\frac{2\pi}{t_0}\mathbb Z\right)=1.
\]
进一步证明，对 \(T>0\)，\(\phi(t+T)=\phi(t)\) 对所有 \(t\in\mathbb R\) 成立，
当且仅当 \(\mathbb P(X\in(2\pi/T)\mathbb Z)=1\)。
\end{optionalexercise}
\begin{solution}
令 \(Z=e^{\mathrm it_0X}\)、\(c=\mathbb EZ\)。展开绝对值平方，得到
\[
 \mathbb E|Z-c|^2=\mathbb E|Z|^2-|c|^2=1-|\phi(t_0)|^2.
\]
若 \(|\phi(t_0)|=1\)，则 \(Z=c\) 几乎必然。
取 \(a\in\mathbb R\) 使 \(c=e^{\mathrm it_0a}\)，便有
\(t_0(X-a)\in2\pi\mathbb Z\) 几乎必然。
反之，这个支集条件使 \(Z\) 几乎必然为同一个模为一的常数，故 \(|\phi(t_0)|=1\)。

若 \(\phi\) 以 \(T\) 为周期，则 \(\phi(T)=\phi(0)=1\)。
以上方差计算取 \(t_0=T\)、\(c=1\)，给出 \(e^{\mathrm iTX}=1\) 几乎必然，
即 \(X\in(2\pi/T)\mathbb Z\) 几乎必然。
反之，该支集条件使
\(e^{\mathrm i(t+T)X}=e^{\mathrm itX}\) 几乎必然，取期望即得周期性。
注意模长等号允许一个平移后的格点集，而特征函数的周期性要求相位恰为 \(1\)。
\end{solution}

% 自拟：以特征函数连续性处理稀有事件的三角阵，区别于固定同分布的CLT。
\begin{optionalexercise}\label{b1:ex:30-poisson-limit}
对每个正整数 \(n\)，设 \(X_{n,1},\ldots,X_{n,n}\) 为某个概率空间上的独立 Bernoulli 变量，
共同参数为 \(p_n\in[0,1]\)。不同的行可以处在不同的概率空间中。
若 \(np_n\to\lambda\in[0,\infty)\)，证明 \(S_n=\sum_{k=1}^nX_{n,k}\) 的分布弱收敛于
\[
 \mu_\lambda=\sum_{j=0}^{\infty}e^{-\lambda}\frac{\lambda^j}{j!}\delta_j,
\]
其中 \(\lambda=0\) 时约定 \(\mu_0=\delta_0\)。
\end{optionalexercise}
\begin{solution}
指数级数说明 \(\mu_\lambda\) 的总质量为一，其特征函数为
\[
 \phi_\lambda(t)
 =e^{-\lambda}\sum_{j=0}^{\infty}\frac{\lambda^je^{\mathrm itj}}{j!}
 =\exp\bigl(\lambda(e^{\mathrm it}-1)\bigr).
\]
独立性给出
\[
 \phi_{S_n}(t)=\bigl(1+p_n(e^{\mathrm it}-1)\bigr)^n.
\]
固定 \(t\)，令 \(z_n=p_n(e^{\mathrm it}-1)\)。由 \(np_n\to\lambda\)，
有 \(nz_n\to\lambda(e^{\mathrm it}-1)\)。
对复数列 \(1+z_n\) 应用\cref{b1:lem:30-complex-power-limit}，得到
\[
 (1+z_n)^n\longrightarrow\exp\bigl(\lambda(e^{\mathrm it}-1)\bigr).
\]
极限是概率测度 \(\mu_\lambda\) 的特征函数，由特征函数连续性定理，得到所需弱收敛。
这里每一行的参数随 \(n\) 改变，和的均值保持有限；没有采用中心极限定理中的平方根归一化。
\end{solution}

% 自拟：经典Rademacher乘积构造，连接两两独立、弱大数律与中心极限失效。
\begin{optionalexercise}\label{b1:ex:30-pairwise-clt-failure}
设 \((\varepsilon_j)_{j\geq1}\) 为同一概率空间上的独立变量列，
\(\mathbb P(\varepsilon_j=1)=\mathbb P(\varepsilon_j=-1)=1/2\)。
把正整数 \(k\) 的二进制展开中取值为 \(1\) 的位置组成的有限非空集记为 \(A(k)\)，最低位的位置为 \(1\)，令
\[
 X_k=\prod_{j\in A(k)}\varepsilon_j,\quad S_n=\sum_{k=1}^nX_k.
\]
证明 \((X_k)\) 两两独立，各项均为取值 \(\pm1\) 的对称分布，且 \(S_n/n\to0\) 依概率成立。
但是，沿 \(n_m=2^m-1\)，有 \(S_{n_m}/\sqrt{n_m}\to0\) 依概率成立，
所以该序列不满足标准正态中心极限结论。
\end{optionalexercise}
\begin{solution}
非空指标集上的独立乘积给出 \(\mathbb EX_k=0\)，且 \(X_k\) 只取 \(\pm1\)，所以各有概率 \(1/2\)。
若 \(k\ne l\)，两个指标集的对称差非空；利用 \(\varepsilon_j^2=1\)，有
\[
 \mathbb E(X_kX_l)
 =\mathbb E\prod_{j\in A(k)\triangle A(l)}\varepsilon_j=0.
\]
对 \(a,b\in\{-1,1\}\)，于是
\[
 \mathbb P(X_k=a,X_l=b)
 =\frac14\mathbb E\bigl((1+aX_k)(1+bX_l)\bigr)=\frac14,
\]
证明了两两独立。方差可加给出 \(\mathbb ES_n=0\)、\(\operatorname{Var}(S_n)=n\)，
故 Chebyshev 估计说明 \(S_n/n\to0\) 依概率成立。

前 \(2^m-1\) 个正整数对应 \(\{1,\ldots,m\}\) 的全部非空子集。展开有限乘积可得
\[
 S_{2^m-1}=\prod_{j=1}^m(1+\varepsilon_j)-1.
\]
除非前 \(m\) 个符号全部为 \(1\)，右侧都等于 \(-1\)；全部为 \(1\) 的事件概率为 \(2^{-m}\)。
因此对任意 \(\delta>0\)，充分大的 \(m\) 满足
\[
 \mathbb P\left(\frac{|S_{2^m-1}|}{\sqrt{2^m-1}}>\delta\right)\leq2^{-m}\longrightarrow0.
\]
该子列的弱极限为 \(\delta_0\)，不可能又是标准正态分布。
两两独立足以消去方差中的交叉项，却不足以把任意有限个特征函数相乘。
\end{solution}

% 自拟：从强大数律估计方差，再用紧性控制随机归一化；不调用未证的统计极限定理。
\begin{optionalexercise}\label{b1:ex:30-estimated-variance}
设 \((\Omega,\mathcal F,\mathbb P)\) 为概率空间，\((X_k)_{k\geq1}\) 为独立同分布的实随机变量，
\(\mathbb EX_1=m\)、\(\operatorname{Var}(X_1)=\sigma^2\in(0,\infty)\)。令
\[
 \overline X_n=\frac1n\sum_{k=1}^nX_k,\quad
 V_n=\frac1n\sum_{k=1}^n(X_k-\overline X_n)^2.
\]
证明 \(V_n\to\sigma^2\) 几乎必然，并证明变量
\[
 T_n=\begin{cases}
 \sqrt n\,(\overline X_n-m)/\sqrt{V_n},&V_n>0,\\
 0,&V_n=0
 \end{cases}
\]
依分布收敛于标准正态分布。证明中须说明随机分母为何不改变极限。
\end{optionalexercise}
\begin{solution}
由恒等式
\(V_n=n^{-1}\sum_{k=1}^nX_k^2-\overline X_n^2\)，
分别对可积变量 \(X_k^2\) 和 \(X_k\) 使用强大数律，得到
\[
 V_n\longrightarrow\mathbb EX_1^2-m^2=\sigma^2\quad\text{几乎必然}.
\]
记 \(Z_n=\sqrt n\,(\overline X_n-m)/\sigma\)。中心极限定理给出
\(Z_n\xrightarrow{d}N(0,1)\)，而且直接由独立性，\(\mathbb EZ_n^2=1\)。
令 \(R_n=\sigma/\sqrt{V_n}\) 于 \(V_n>0\)，在其余点取零，则 \(T_n=R_nZ_n\)。
因为 \(V_n\to\sigma^2>0\) 几乎必然，有 \(R_n\to1\) 依概率。
对任意 \(M>0\)、\(\delta>0\)，
\[
 \mathbb P(|T_n-Z_n|>\delta)
 \leq\mathbb P(|Z_n|>M)+\mathbb P(|R_n-1|>\delta/M)
 \leq M^{-2}+\mathbb P(|R_n-1|>\delta/M).
\]
先令 \(n\to\infty\)，再令 \(M\to\infty\)，得到 \(T_n-Z_n\to0\) 依概率。

最后不预先引用随机归一化定理，而直接比较分布函数。
对每个 \(a\in\mathbb R\) 和 \(\delta>0\)，有
\[
 \begin{split}
 \mathbb P(Z_n\leq a-\delta)-\mathbb P(|T_n-Z_n|>\delta)
 &\leq\mathbb P(T_n\leq a),\\
 \mathbb P(T_n\leq a)
 &\leq\mathbb P(Z_n\leq a+\delta)+\mathbb P(|T_n-Z_n|>\delta).
 \end{split}
\]
令 \(n\to\infty\)，上下界分别趋于 \(\Phi(a-\delta)\)、\(\Phi(a+\delta)\)。
标准正态分布函数 \(\Phi\) 连续，再令 \(\delta\downarrow0\)，便得
\(\mathbb P(T_n\leq a)\to\Phi(a)\)。由分布函数判据，证明了所需依分布收敛。
\end{solution}
